% BEGIN INCLUDED SGA7I_English_current_buildable_source_20260823/components/178_expose_IX_sections_13_1_and_13_2.tex
\medskip
\noindent
13. \underline{与超越理论的联系：复解析情形}.

关于解析空间的一般概念，参见 [2]。

\medskip
\noindent
\SGARefTarget{sga7.sga7i.eix.n13-1}{13.1}. \underline{$S$ 上解析群的回顾。} 设 $S$ 是解析空间，$A$ 是 $S$ 上的解析空间，带有 $S$ 上的交换群结构（即 $A$ 是 $S$ 上解析空间范畴中的交换群）。假设 $A$ 在 $S$ 上\underline{光滑}，并以

\begin{equation}\tag{\SGARefTarget{sga7.sga7i.eix.n13-1-1}{13.1.1}}
 \underline E=\underline{\operatorname{Lie}}(A/S)
\end{equation}

\noindent
表示把 $A$ 在 $S$ 上的相对切层沿零截面拉回得到的 $S$ 上局部自由模。在此承认指数同态

\begin{equation}\tag{\SGARefTarget{sga7.sga7i.eix.n13-1-2}{13.1.2}}
 \exp:E\longrightarrow A\qquad ,
\end{equation}

\noindent
的定义，其中 $E$ 是 $S$ 上的解析向量丛，其全纯截面模为 $\underline E$。在点 $s\in S$ 的邻域中，先用特征零形式群理论的熟知方法，在 $E$ 的零截面附近定义此同态，再显然地利用乘法相容性将其延拓到整个 $E$。它是 $S$ 上相对解析群的同态；它是\underline{平展的}，其像为 $A$ 的开子空间 $A^o$，后者是 $A$ 在 $S$ 上各纤维的单位分支之并。若 $R$ 是作为 $\exp$ 的核的解析子空间，则有相对解析群的正合列

\begin{equation}\tag{\SGARefTarget{sga7.sga7i.eix.n13-1-3}{13.1.3}}
 0\longrightarrow R\longrightarrow E\longrightarrow A^o
 \longrightarrow0\qquad .
\end{equation}

\noindent
这里态射 $E\longrightarrow A^o$ 满射且平展，故 $R$ 是 $S$ 上\underline{平展}的相对解析群。它是 $E$ 的解析子空间，而且它在 $E$ 中闭，当且仅当 $A^o$ 在 $S$ 上分离，即当且仅当 $A^o$ 的零截面在 $A^o$ 中闭。

\medskip
\noindent
\SGARefTarget{sga7.sga7i.eix.n13-1-4}{13.1.4}. 当然，正合列（\SGARefLink{sga7.sga7i.eix.n13-1-3}{13.1.3}）\underline{对 $A$ 具有函子性}，并与每个基变换 $S'\longrightarrow S$ 可交换。由此在 $S$ 上光滑且纤维连通的相对解析群范畴，与由 $S$ 上局部自由向量丛 $E$ 及 $E$ 的一个在 $S$ 上平展的解析子群 $R$ 所成的二元组 $(E,R)$ 的范畴之间，得到一个与基变换相容的范畴等价。

\medskip
\noindent
\SGARefTarget{sga7.sga7i.eix.n13-1-5}{13.1.5}. 当 $A$ 是 $S$ 上的“\underline{类阿贝尔}（abeloid）”群，即它在 $S$ 上光滑、紧合且分离，并且纤维连通——故特别有 $A=A^o$——时，$R$ 是 $S$ 上局部常值的群，在每个点 $s$ 的秩（作为 $\mathbb Z$-模）等于 $2\operatorname{rank}_s(E)$，反之亦然。由此得到对 $S$ 上类阿贝尔空间的熟知解释：它由 $S$ 上的一个“\underline{局部系统}” $R$，以及在 $R$ 定义的实解析丛 $R\otimes_{\mathbb Z}\mathbb R$ 上的一个\underline{复}解析向量丛结构给出。

\medskip
\noindent
\SGARefTarget{sga7.sga7i.eix.n13-2}{13.2}. \underline{沿一个纤维的形式完备化。} 现在固定点 $s\in S$，在 $s$ 附近局部地工作。对每个整数 $n\geqslant0$，以 $S_n$ 表示 $s$ 在 $S$ 中的第 $n$ 阶无穷小邻域，并对 $S$ 上的每个解析空间 $X$，同样置 $X_n=X\times_S S_n$。当 $n$ 变化时，$X_n$ 所成的系统构成 $S$ 上解析空间的归纳系统，记为 $\widehat X$，称为\underline{$X$ 沿 $X_s$ 的形式完备化}。它显然对 $X$ 具有函子性，其形成与有限投射极限可交换；所以 $X\longmapsto\widehat X$ 将群映为群。注意 $\widehat X$ 实际上只依赖于 $X$ 在 $s$ 在 $S$ 中任意一个开邻域 $U$ 上的限制。

\medskip
\noindent
\underline{引理 \SGARefTarget{sga7.sga7i.eix.n13-2-1}{13.2.1}.} \underline{数据如 \SGARefLink{sga7.sga7i.eix.n13-1}{13.1}，并假设下列条件成立：} a) $A=A^o$；b) $R$ \underline{在 $s$ 的一个邻域中为常值的}；以及 c) $R_s$ \underline{张成复向量空间 $E_s$}。\underline{设 $A'$ 是 $S$ 上的另一个光滑解析群，并考虑映射}

\begin{equation*}
 u\longmapsto\widehat u:
 \operatorname{Hom}_{S\text{-}\mathrm{gr}}(A,A')
 \longrightarrow
 \operatorname{Hom}_{\widehat S\text{-}\mathrm{gr}}
   (\widehat A,\widehat {A}')\quad ,
\end{equation*}

\noindent
\underline{其中记号如 \SGARefLink{sga7.sga7i.eix.n13-2}{13.2}；再考虑限制到 $s$ 的可变开邻域 $U$ 并取归纳极限所得到的类似映射：}

\begin{equation}\tag{\SGARefTarget{sga7.sga7i.eix.n13-2-1-1}{13.2.1.1}}
 u\longmapsto\widehat u:
 \operatorname{Hom}_{(S,s)\text{-}\mathrm{gr}}(A,A')
 \longrightarrow
 \operatorname{Hom}_{\widehat S\text{-}\mathrm{gr}}
   (\widehat A,\widehat {A}')\quad .
\end{equation}

\noindent
（\underline{这里 $(S,s)$ 表示解析空间 $S$ 在 $s$ 处的芽。}）\underline{则后一个映射为双射。}

事实上，考虑正合列

\begin{equation*}
 0\longrightarrow R'\longrightarrow E'\longrightarrow {A'}^o
 \longrightarrow C
\end{equation*}

\noindent
它类似于（\SGARefLink{sga7.sga7i.eix.n13-1-3}{13.1.3}），由 $A'$ 定义。由于同态 $A\longrightarrow A'$ 必然经过 ${A'}^o$ 分解，而且在基变换 $U\longrightarrow S$ 和 $S_n\longrightarrow S$ 后也如此，可假设 $A'={A'}^o$。由 \SGARefLink{sga7.sga7i.eix.n13-1-4}{13.1.4} 中的对应，只须证明：给出把 $R$ 映入 $R'$ 的向量丛同态 $E\longrightarrow E'$ 的一个芽，等价于给出把 $R_n$ 映入 $R'_n$ 的一族相容同态 $E_n\longrightarrow E'_n$。

若 $\underline{\mathcal O}_s$ 表示 $S$ 在 $s$ 处的局部环，给出一族相容同态 $E_n\longrightarrow E'_n$，等价于给出一个 $\widehat{\underline{\mathcal O}}_s$-模同态

\begin{equation*}
 \widehat F:\widehat{\underline E}_s
 \longrightarrow\widehat{\underline E}'_s,
\end{equation*}

\noindent
其中 $\underline E_s$ 和 $\underline E'_s$ 是由分别对应于 $E$ 和 $E'$ 的 $\underline{\mathcal O}_s$-模 $\underline E$ 和 $\underline E'$ 在 $s$ 处得到的自由 $\underline{\mathcal O}_s$-模。另一方面，向量丛同态 $E\longrightarrow E'$ 的一个芽，等价于给出一个 $\underline{\mathcal O}_s$-模同态

\begin{equation*}
 F:\underline E_s\longrightarrow\underline E'_s\quad .
\end{equation*}

\noindent
这已说明（\SGARefLink{sga7.sga7i.eix.n13-2-1-1}{13.2.1.1}）为单射。为证满射性，只须证明：若 $\widehat F$ 为如上的同态，且对每个 $n$，所诱导的同态 $F_n:E_n\longrightarrow E'_n$ 把 $R_n$ 映入 $R'_n$，则 $\widehat F$ 来自某个 $F$——即它把 $\underline E_s$ 映入 $\underline E'_s$——并且相应的同态芽 $E\longrightarrow E'$ 把 $R$ 映入 $R'$。

由假设 c)，$R_s=R_o$ 张成 $E_s$；由 Nakayama 引理，群 $R_o$（可等同于 $\widehat R$ 在 $\widehat S$ 上的截面群，也可等同于 $R$ 在 $S$ 上截面芽的群）生成 $\widehat E_s$。因此，为证 $\widehat F(\underline E_s)\subset\underline E'_s$，只须证 $\widehat F(R_o)\subset\underline E'_s$。但 $\widehat F(R_o)\subset R'_o$，而 $R'_o$ 本身可解释为 $R'\subset E'$ 在 $s$ 附近的截面芽所成的群。这给出所述包含关系，从而得到 $F$ 的存在性。此外可见 $F(R_o)\subset R'_o$，其中 $R_o$ 和 $R'_o$ 分别视为 $E$ 和 $E'$ 的截面芽的群。由 b)，$R_o$ 在 $s$ 的一个邻域中生成 $R$；故在 $s$ 的一个邻域中有 $F(R)\subset R'$，证明完毕。

\medskip
\noindent
\underline{推论 \SGARefTarget{sga7.sga7i.eix.n13-2-2}{13.2.2}.} \underline{若 $A$ 满足 \SGARefLink{sga7.sga7i.eix.n13-2-1}{13.2.1} 的条件 a)、b) 和 c)，则在 $s$ 的一个邻域中，它由 $\widehat S$-群 $\widehat A$ 在唯一同构意义下确定。}\footnote{印本此处指向 \SGARefLink{sga7.sga7i.eix.n13-2-2}{13.2.2} 自身的条件；这三个条件引入于引理 \SGARefLink{sga7.sga7i.eix.n13-2-1}{13.2.1}。}

% END INCLUDED SGA7I_English_current_buildable_source_20260823/components/178_expose_IX_sections_13_1_and_13_2.tex

% BEGIN INCLUDED SGA7I_English_current_buildable_source_20260823/components/179_expose_IX_section_13_3.tex
\medskip
\noindent
\SGARefTarget{sga7.sga7i.eix.n13-3}{13.3}. \underline{解析 Raynaud 扩张}.

现在可以在解析几何中仿照第7节给出的 Raynaud 扩张构造。假设 $A$ 如 \SGARefLink{sga7.sga7i.eix.n13-1}{13.1}，且 $A=A^o$，\footnote{印本此处为 $A=A_o$。所需并在下文使用的假设是 $A=A^o$；比较 \SGARefLink{sga7.sga7i.eix.n13-1}{13.1} 和引理 \SGARefLink{sga7.sga7i.eix.n13-2-1}{13.2.1}。}并考虑 $R_o=R_s$，它是有限型自由 $\mathbb Z$-模。将 $S$ 限制到 $s$ 的一个充分小的开邻域后，把 $R_o$ 解释为 $R$ 在 $s$ 附近的截面芽所成的群，得到同态

\begin{equation}\tag{\SGARefTarget{sga7.sga7i.eix.n13-3-1}{13.3.1}}
 R_{oS}\longrightarrow R\subset E,
\end{equation}

\noindent
其中 $R_{oS}$ 表示 $S$ 上取值为 $R_o$ 的常值解析群。选取 $s$ 的邻域足够小，使 $\underline{\mathcal O}_S$ 的每个在 $s$ 处的芽为零的截面本身都为零，则可见上述同态为单射。因此其像是 $R$ 的常值解析子群；显然它是 $R$ 的最大常值子群，称为 $R$ 的\underline{不变部分}，记为 $R^f$：

\begin{equation}\tag{\SGARefTarget{sga7.sga7i.eix.n13-3-2}{13.3.2}}
 R_{oS}\simeq R^f\hookrightarrow R,
 \qquad\text{诱导同构 }R_s^f\simeq R_s.
\end{equation}

\noindent
于是可以考虑商群

\begin{equation}\tag{\SGARefTarget{sga7.sga7i.eix.n13-3-3}{13.3.3}}
 A^\natural=E/R^f,
\end{equation}

\noindent
并由（\SGARefLink{sga7.sga7i.eix.n13-1-3}{13.1.3}）得到典范正合列

\begin{equation}\tag{\SGARefTarget{sga7.sga7i.eix.n13-3-4}{13.3.4}}
 0\longrightarrow R/R^f\stackrel{i}{\longrightarrow}A^\natural
 \longrightarrow A\longrightarrow0.
\end{equation}

\noindent
注意，将 $S$ 限制到 $s$ 的邻域后，$R^f$ 总在 $E$ 中闭，故 $A^\natural$ 在 $S$ 上分离。还须注意同态 $A^\natural\longrightarrow A$ 平展，因而诱导同构

\begin{equation}\tag{\SGARefTarget{sga7.sga7i.eix.n13-3-5}{13.3.5}}
 \underline{\operatorname{Lie}}(A^\natural)
 \stackrel{\sim}{\longrightarrow}
 \underline{\operatorname{Lie}}(A)=\underline E.
\end{equation}

\noindent
此外，对每个 $n\geqslant0$，典范同态 $A^\natural\longrightarrow A$ 诱导同构 $A_n^\natural\simeq A_n$（因为 $R_n^f\stackrel{\sim}{\longrightarrow}R_n$），从而诱导同构

\begin{equation}\tag{\SGARefTarget{sga7.sga7i.eix.n13-3-6}{13.3.6}}
 \widehat{A^\natural}\stackrel{\sim}{\longrightarrow}\widehat A.
\end{equation}

\noindent
最后这个事实说明，$A^\natural$ 的确起着第7节中同样记为 $A^\natural$ 的群的作用。这里的构造对纤维 $A_s=A_o$ 未作“类阿贝尔群对环面的扩张”这类假设。然而，若要在 $A^\natural$ 上恢复重要的扩张结构（\SGARefLink{sga7.sga7i.eix.n7-2-3}{7.2.3}），这种假设仍将是必要的。相比之下，（\SGARefLink{sga7.sga7i.eix.n13-3-4}{13.3.4}）的同态 $A^\natural\longrightarrow A$ 本质上具有超越性质，在第7节的纯代数理论中没有对应物。第14节将再讨论（\SGARefLink{sga7.sga7i.eix.n13-3-4}{13.3.4}）中同态 $i:R/R^f\longrightarrow A^\natural$ 的纯代数解释问题。

先假设给定一个复子环面

\begin{equation}\tag{\SGARefTarget{sga7.sga7i.eix.n13-3-7}{13.3.7}}
 T_o\subset A_o,
\end{equation}

\noindent
其中“\underline{复环面}”是指同构于 $\mathbb C^*$ 的某个笛卡儿幂的解析群。设

\begin{equation}\tag{\SGARefTarget{sga7.sga7i.eix.n13-3-8}{13.3.8}}
 T=T_o\times S
\end{equation}

\noindent
为 $S$ 上取值为 $T_o$ 的常值环面。按 \SGARefLink{sga7.sga7i.eix.n5-1}{5.1} 的方法，得到解析群的归纳系统的唯一同态（实际上是浸入）

\begin{equation*}
 \widehat T\longrightarrow\widehat{A^\natural},
\end{equation*}

\noindent
它在 $S_o$ 上约化为（\SGARefLink{sga7.sga7i.eix.n13-3-7}{13.3.7}）。由 \SGARefLink{sga7.sga7i.eix.n13-2-2}{13.2.2}，此同态来自唯一的同态芽

\begin{equation}\tag{\SGARefTarget{sga7.sga7i.eix.n13-3-9}{13.3.9}}
 T\longrightarrow A^\natural,
\end{equation}

\noindent
在 $s$ 附近把 $S$ 取得足够小后，可假设它定义在整个 $S$ 上。相应的同态

\begin{equation*}
 \underline{\operatorname{Lie}}(T)\longrightarrow
 \underline{\operatorname{Lie}}(A^\natural)\simeq\underline E
\end{equation*}

\noindent
在每个 $S_n$ 上单射，故在 $\widehat S$ 上单射，从而由平坦下降，在 $\underline{\mathcal O}_s$ 上单射。因此，$s$ 在 $S$ 中有一个邻域（不妨取为整个 $S$），使 $\underline{\operatorname{Lie}}(T)$ 等同于 $\underline E$ 的一个局部直和项 $\underline E^t$；以 $E^t$ 表示 $E$ 的相应向量子丛：

\begin{equation}\tag{\SGARefTarget{sga7.sga7i.eix.n13-3-10}{13.3.10}}
 \underline{\operatorname{Lie}}(T)\xrightarrow{\ \sim\ }
 \underline E^t\subset\underline E\simeq
 \underline{\operatorname{Lie}}(A^\natural)\simeq
 \underline{\operatorname{Lie}}(A),
 \qquad E^t\subset E.
\end{equation}

\noindent
类似地，令

\begin{equation}\tag{\SGARefTarget{sga7.sga7i.eix.n13-3-11}{13.3.11}}
 R_o^t\subset R_o,
 \qquad R^t=R_{oS}^t\subset R_{oS}\simeq R^f
\end{equation}

\noindent
分别表示 $A_o$ 的子环面 $T_o$ 所定义的 $R_o$ 的子格，以及 $S$ 上相应的常值解析群。于是，同态（\SGARefLink{sga7.sga7i.eix.n13-3-9}{13.3.9}）仅对应于同态

\begin{equation*}
 E^t/R_{oS}^t=E^t/R^t\longrightarrow E/R_{oS}=E/R^f
\end{equation*}

\noindent
它由下列包含诱导：

\begin{equation*}
 E^t\hookrightarrow E,
 \qquad R^t\hookrightarrow R^f.
\end{equation*}

\noindent
利用（\SGARefLink{sga7.sga7i.eix.n13-3-9}{13.3.9}）在 $s$ 上的纤维上诱导浸入这一事实，得到同态

\begin{equation*}
 R_o/R_o^t\longrightarrow E_o/E_o^t=(E/E^t)_o
\end{equation*}

\noindent
为单射。立即得到

\begin{equation*}
 R_o/R_o^t\longrightarrow(E/E^t)_{s'}
\end{equation*}

\noindent
对于靠近 $s$ 的 $s'$ 仍为单射。因此，在 $s$ 附近缩小 $S$ 后，可假设（\SGARefLink{sga7.sga7i.eix.n13-3-9}{13.3.9}）\underline{为闭浸入}。

因此引入

\begin{equation}\tag{\SGARefTarget{sga7.sga7i.eix.n13-3-12}{13.3.12}}
 B=A^\natural/T\simeq(E/E^t)/(R^f/R^t),
\end{equation}

\noindent
从而 $A^\natural$ 获得扩张结构

\begin{equation}\tag{\SGARefTarget{sga7.sga7i.eix.n13-3-13}{13.3.13}}
 0\longrightarrow T\longrightarrow A^\natural
 \longrightarrow B\longrightarrow0,
\end{equation}

\noindent
其中 $B$ 是 $S$ 上的相对解析群，在 $S$ 上光滑且分离，并且纤维连通。因此由构造，有

\begin{equation}\tag{\SGARefTarget{sga7.sga7i.eix.n13-3-14}{13.3.14}}
 \underline{\operatorname{Lie}}(B)\simeq
 \underline E/\underline E^t,
\end{equation}

\noindent
而 $B$ 的指数正合列为

\begin{equation}\tag{\SGARefTarget{sga7.sga7i.eix.n13-3-15}{13.3.15}}
 0\longrightarrow R^f/R^t\longrightarrow E/E^t
 \longrightarrow B\longrightarrow0.
\end{equation}

\noindent
注意

\begin{equation}\tag{\SGARefTarget{sga7.sga7i.eix.n13-3-16}{13.3.16}}
 R^f/R^t=(R_o/R_o^t)_S
\end{equation}

\noindent
是 $S$ 上的常值解析群。由 \SGARefLink{sga7.sga7i.eix.n13-1-5}{13.1.5}，$B$ 在 $s$ 的一个邻域中是 $S$ 上的类阿贝尔群，当且仅当它在 $s$ 处的纤维

\begin{equation*}
 B_o=A_o/T_o,
\end{equation*}

\noindent
是类阿贝尔解析群；此时 $A_o$ 表为此类群 $B_o$ 对环面 $T_o$ 的扩张。因此，\underline{把 $A_o$ 写为类阿贝尔群对复环面的扩张 $(*)$ 的每一种方式，都在 $s$ 的一个邻域中给 $A$ 赋予唯一的扩张结构}（\SGARefLink{sga7.sga7i.eix.n13-3-13}{13.3.13}），\underline{即类阿贝尔群 $B$ 对环面 $T$ 的扩张，并在 $s$ 处诱导给定的扩张结构}。

% END INCLUDED SGA7I_English_current_buildable_source_20260823/components/179_expose_IX_section_13_3.tex


% BEGIN INCLUDED SGA7I_English_current_buildable_source_20260823/components/180_expose_IX_section_13_4.tex
\medskip
\noindent
\SGARefTarget{sga7.sga7i.eix.n13-4}{13.4}. \underline{用 Deligne 的“混合 Hodge 结构”解释.}

先回到 \SGARefLink{sga7.sga7i.eix.n13-1-4}{13.1.4} 中的对应，将它略作重述，使之更适于引入当前情形有关的“Hodge 结构”。为此，设 $\underline R$ 是 $R$ 在 $S$ 上的截面层，并引入 $\underline{\mathcal O}_S$-模

\begin{equation}\tag{\SGARefTarget{sga7.sga7i.eix.n13-4-1}{13.4.1}}
 \mathcal E(A)=\mathcal E=
 \underline R\otimes_{\mathbb Z}\underline{\mathcal O}_S\qquad .
\end{equation}

\noindent
当 $R$ 局部常值时，它是有限型局部自由的。（注意，当 $R$ 不局部常值时，$\mathcal E$ 未必凝聚。）于是（\SGARefLink{sga7.sga7i.eix.n13-1-3}{13.1.3}）中的同态 $\underline R\longrightarrow\underline E$ 定义一个典范同态

\begin{equation}\tag{\SGARefTarget{sga7.sga7i.eix.n13-4-2}{13.4.2}}
 \mathcal E=\underline R\otimes_{\mathbb Z}
 \underline{\mathcal O}_S\longrightarrow\underline E\qquad .
\end{equation}

\noindent
将其核记为 $\underline F(A)$，或简记为 $\underline F$；因此有正合列

\begin{equation}\tag{\SGARefTarget{sga7.sga7i.eix.n13-4-3}{13.4.3}}
 0\longrightarrow\underline F\longrightarrow\mathcal E
 \longrightarrow\underline E\qquad .
\end{equation}

\noindent
最有意义的情形是（\SGARefLink{sga7.sga7i.eix.n13-4-2}{13.4.2}）为满态射，即对每个 $s\in S$，$R_s$ 都张成复向量空间 $E_s$。此时有正合列

\begin{equation}\tag{\SGARefTarget{sga7.sga7i.eix.n13-4-4}{13.4.4}}
 0\longrightarrow\underline F\longrightarrow\mathcal E
 \longrightarrow\underline E\longrightarrow0\qquad .
\end{equation}

\noindent
若 $R$ 局部常值，这是 $S$ 上局部自由模的短正合列。此时给出二元组 $(E,R)$，等价于给出二元组 $(R,\underline F)$，其中 $R$ 是 $S$ 上局部常值的解析群，其纤维为有限型自由 $\mathbb Z$-模，$\underline F$ 是 $\mathcal E=\underline R\otimes_{\mathbb Z}\underline{\mathcal O}_S$ 的局部直和项子模，并另须满足

\begin{equation}\tag{\SGARefTarget{sga7.sga7i.eix.n13-4-5}{13.4.5}}
 \underline F\cap\underline R=0\qquad .
\end{equation}

\noindent
从 Hodge 结构的角度，应把 $\mathcal E$ 看作带有“\underline{Hodge 滤过}”

\begin{equation}\tag{\SGARefTarget{sga7.sga7i.eix.n13-4-6}{13.4.6}}
 \operatorname{Fil}^{1}(\mathcal E)=\mathcal E,
 \qquad \operatorname{Fil}^{0}(\mathcal E)=\underline F,
 \qquad \operatorname{Fil}^{-1}(\mathcal E)=0\qquad .
\end{equation}

\noindent
注意，正合列（\SGARefLink{sga7.sga7i.eix.n13-4-3}{13.4.3}）仍对 $S$ 上的解析群 $A$ 具有函子性，并与任意基变换可交换。

现在回到 \SGARefLink{sga7.sga7i.eix.n13-3}{13.3} 的条件，置

\begin{equation}\tag{\SGARefTarget{sga7.sga7i.eix.n13-4-7}{13.4.7}}
 \mathcal E^f=\underline R^f\otimes_{\mathbb Z}
 \underline{\mathcal O}_S,
 \qquad
 \mathcal E^t=\underline R^t\otimes_{\mathbb Z}
 \underline{\mathcal O}_S\qquad .
\end{equation}

\noindent
这些模定义在 $s$ 的某个适当邻域上，不妨假设此邻域为整个 $S$；它们定义 $\mathcal E$ 的一个称为\underline{权滤过}的滤过，

\begin{equation}\tag{\SGARefTarget{sga7.sga7i.eix.n13-4-8}{13.4.8}}
 \mathcal E\supset\mathcal E^f\supset\mathcal E^t\supset0\qquad .
\end{equation}

\noindent
这里应把 $\mathcal E^t$、$\mathcal E^f$ 和 $\mathcal E$ 分别看作具有权 $-2$、$-1$ 和 $0$；因此这是递增滤过。与 $R^t$ 一样，模 $\mathcal E^t$ 依赖于 $A_o=A_s$ 的子环面 $T_o$ 的选取。由于 $R^f$ 和 $R^t$ 局部常值，\underline{$\mathcal E^f$ 和 $\mathcal E^t$ 局部自由}；还可将它们解释为

\begin{equation}\tag{\SGARefTarget{sga7.sga7i.eix.n13-4-9}{13.4.9}}
 \mathcal E^f=\mathcal E(A^\natural),
 \qquad \mathcal E^t=\mathcal E(T)\qquad .
\end{equation}

\noindent
将（\SGARefLink{sga7.sga7i.eix.n13-4-3}{13.4.3}）对可变 $A$ 的函子性应用于同态

\begin{equation*}
 T\longrightarrow A^\natural\longrightarrow A
\end{equation*}

\noindent
立即得到，$\underline F(A^\natural)$ 和 $\underline F(T)$ 分别等同于

\begin{equation}\tag{\SGARefTarget{sga7.sga7i.eix.n13-4-10}{13.4.10}}
 \underline F(A^\natural)\simeq
 \underline F^f\stackrel{\mathrm{def}}{=}
 \underline F\cap\mathcal E^f,
 \qquad
 \underline F(T)=\underline F^t=
 \underline F\cap\mathcal E^t=0\qquad .
\end{equation}

\noindent
最后的等式直接来自环面 $T$ 满足 $\underline F(T)=0$ 这一事实：此时（\SGARefLink{sga7.sga7i.eix.n13-4-3}{13.4.3}）型的正合列是 $0\longrightarrow0\longrightarrow\mathcal E(T)\longrightarrow\underline{\operatorname{Lie}}(T)\longrightarrow0$。同时看到，复合同态

\begin{equation*}
 \mathcal E^t\longrightarrow\mathcal E\longrightarrow\underline E
\end{equation*}

\noindent
诱导同构

\begin{equation}\tag{\SGARefTarget{sga7.sga7i.eix.n13-4-11}{13.4.11}}
 \mathcal E^t\stackrel{\sim}{\longrightarrow}\underline E^t
 \simeq\underline{\operatorname{Lie}}(T)\qquad .
\end{equation}

\noindent
\underline{现在假设 $R_s$ 张成复向量空间 $E_s$}，特别地，在 \SGARefLink{sga7.sga7i.eix.n13-3-16}{13.3.16} 的有利情形中，即 $A_s$ 为类阿贝尔群对复环面 $T_o$ 的扩张时如此。将 $S$ 限制到 $s$ 的一个适当邻域后，以 $A^\natural$ 代替 $A$ 所得的（\SGARefLink{sga7.sga7i.eix.n13-4-2}{13.4.2}）型典范同态

\begin{equation*}
 \mathcal E^f=\mathcal E(A^\natural)
 \longrightarrow\underline{\operatorname{Lie}}(A^\natural)
 \simeq\underline E
\end{equation*}

\noindent
为满态射；于是（\SGARefLink{sga7.sga7i.eix.n13-4-2}{13.4.2}）更是如此。因此有正合列（\SGARefLink{sga7.sga7i.eix.n13-4-4}{13.4.4}）以及

\begin{equation}\tag{\SGARefTarget{sga7.sga7i.eix.n13-4-12}{13.4.12}}
 0\longrightarrow\underline F^f\longrightarrow\mathcal E^f
 \longrightarrow\underline E\longrightarrow0\qquad .
\end{equation}

\noindent
特别地，

\begin{equation}\tag{\SGARefTarget{sga7.sga7i.eix.n13-4-13}{13.4.13}}
 \underline F+\mathcal E^f=\mathcal E\qquad ,
\end{equation}

\noindent
因为 $\mathcal E/\underline F\simeq\underline E$。此外，由于 $B=A^\natural/T$ 是 $A^\natural$ 的商，\footnote{印本此处说 $B$ 是 $A$ 的商；所显示的定义 $B=A^\natural/T$ 给出了所指的群。}它满足类似的满射性条件，因而给出（\SGARefLink{sga7.sga7i.eix.n13-4-4}{13.4.4}）型的正合列：

\begin{equation}\tag{\SGARefTarget{sga7.sga7i.eix.n13-4-14}{13.4.14}}
 0\longrightarrow\underline F(B)\longrightarrow\mathcal E(B)
 \longrightarrow\underline{\operatorname{Lie}}(B)
 \longrightarrow0\qquad .
\end{equation}

\noindent
用已经引入的模解释此列是直接的。事实上，由（\SGARefLink{sga7.sga7i.eix.n13-3-12}{13.3.12}），

\begin{equation}\tag{\SGARefTarget{sga7.sga7i.eix.n13-4-15}{13.4.15}}
 \left\{
 \begin{array}{l}
  \mathcal E(B)\simeq
  (\underline R^f/\underline R^t)\otimes_{\mathbb Z}
  \underline{\mathcal O}_S
  \simeq\mathcal E^f/\mathcal E^t,\\[2ex]
  \underline{\operatorname{Lie}}(B)
  \simeq\underline E/\underline E^t.
 \end{array}
 \right.
\end{equation}

\noindent
由（\SGARefLink{sga7.sga7i.eix.n13-4-11}{13.4.11}）及（\SGARefLink{sga7.sga7i.eix.n13-4-10}{13.4.10}）中的关系 $\underline F\cap\mathcal E^t=0$，得到 $\underline F(B)=\ker(\mathcal E^f/\mathcal E^t\longrightarrow\underline E/\underline E^t)$ 通过下列同态同构于 $\underline F^f$：

\begin{equation}\tag{\SGARefTarget{sga7.sga7i.eix.n13-4-16}{13.4.16}}
 \underline F^f\simeq\underline F(A^\natural)
 \stackrel{\sim}{\longrightarrow}\underline F(B)\qquad .
\end{equation}

\noindent
在 \SGARefLink{sga7.sga7i.eix.n13-3-16}{13.3.16} 的有利情形中，即 $A_s$ 为类阿贝尔解析群 $B_o$ 对环面 $T_o$ 的扩张时，总结刚得到的权滤过（\SGARefLink{sga7.sga7i.eix.n13-4-8}{13.4.8}）的性质，以及它与由子模 $\underline F$ 定义的 $\mathcal E$ 的 Hodge 滤过（\SGARefLink{sga7.sga7i.eix.n13-4-6}{13.4.6}）的关系。我们用此环面定义 $R^t$，从而 $B$ 是 $S$ 上的类阿贝尔解析群。

\medskip
\noindent
\underline{命题 \SGARefTarget{sga7.sga7i.eix.n13-4-17}{13.4.17}.} \underline{在上述条件下，$\mathcal E=\underline R\otimes_{\mathbb Z}\underline{\mathcal O}_S$ 上的 Hodge 滤过和权滤过}，

\begin{equation}\tag{\SGARefTarget{sga7.sga7i.eix.n13-4-17-1}{13.4.17.1}}
 \mathcal E\supset\underline F,
 \qquad
 \mathcal E\supset\mathcal E^f\supset\mathcal E^t\supset0,
\end{equation}

\noindent
\underline{满足以下条件：}

\smallskip
\noindent
a) \underline{权滤过由 $\underline R$ 上的类似滤过，对 $\mathbb Z$ 与 $\underline{\mathcal O}_S$ 作张量积得到：}

\begin{equation*}
 \underline R\supset\underline R^f\supset\underline R^t\supset0,
\end{equation*}

\noindent
\underline{其中 $\underline R^f$ 和 $\underline R^t$ 是 $S$ 上局部常值的解析群，取值为有限型自由 $\mathbb Z$-模}，即 $R_o^f=R_s=\pi_1(A_s)$ 和 $R_o^t=\pi_1(T_o)$。

\smallskip
\noindent
b) \underline{有关系}

\begin{equation}\tag{*}
 \underline F\cap\mathcal E^t=0,
 \qquad \underline F+\mathcal E^f=\mathcal E\qquad .
\end{equation}

\noindent
\underline{换言之，由 Hodge 滤过在 $\mathcal E^t$ 上诱导的滤过}，采用（\SGARefLink{sga7.sga7i.eix.n13-4-6}{13.4.6}）的次数约定，\underline{集中于 $-1$ 次}，而 Hodge 滤过在 $\mathcal E/\mathcal E^f$ 上诱导的滤过\underline{集中于 $0$ 次}。

\smallskip
\noindent
c) \underline{在}

\begin{equation*}
 \mathcal E^{\mathrm{ab}}
 \stackrel{\mathrm{def}}{=}\mathcal E^f/\mathcal E^t
 \simeq\underline R^{\mathrm{ab}}\otimes_{\mathbb Z}
 \underline{\mathcal O}_S,
 \qquad
 \underline R^{\mathrm{ab}}=\underline R^f/\underline R^t
 \quad\text{（参见 a)）},
\end{equation*}

\noindent
\underline{上考虑由 $\mathcal E$ 上的 Hodge 滤过诱导的滤过}。因此它由 $\mathcal E^{\mathrm{ab}}$ 的子模 $\underline F^{\mathrm{ab}}$ 定义，后者是下列态射的像：

\begin{equation*}
 \underline F^f\stackrel{\mathrm{def}}{=}
 \underline F\cap\mathcal E^f\longrightarrow\mathcal E^{\mathrm{ab}}.
\end{equation*}

\noindent
由第一个关系 $(*)$，此箭头是单态射。\underline{于是二元组}

\begin{equation*}
 \bigl(\underline R^{\mathrm{ab}},
 \underline F^{\mathrm{ab}}\subset
 \underline R^{\mathrm{ab}}\otimes_{\mathbb Z}
 \underline{\mathcal O}_S=\mathcal E^{\mathrm{ab}}\bigr)
\end{equation*}

\noindent
\underline{同构于 $S$ 上一个类阿贝尔群所联系的二元组} $(\underline R(B),\underline F(B))$，其含义如（\SGARefLink{sga7.sga7i.eix.n13-4-4}{13.4.4}）后所述。因此

\begin{equation*}
 B\simeq\mathcal E^{\mathrm{ab}}/
 (F^{\mathrm{ab}}+R^{\mathrm{ab}}),
\end{equation*}

\noindent
其中不带下划线的符号表示 $S$ 上相应的解析群。

\medskip
\noindent
\underline{注 \SGARefTarget{sga7.sga7i.eix.n13-4-18}{13.4.18}.} 用 Deligne [4] 的用语，把前一个命题的结论限制到 $S$ 上使 $R$ 局部常值的开子集 $U$ 后，就是说：$R$ 配备 $\mathcal E=\underline R\otimes_{\mathbb Z}\underline{\mathcal O}_S$ 上的两个滤过（\SGARefLink{sga7.sga7i.eix.n13-4-17-1}{13.4.17.1}），是“$U$ 上的混合 Hodge 结构解析族”。刚才描述的情形是 Deligne 关于更高维 N\'eron 理论的猜想的模型，这些猜想述于 [9，9.8--9.13]。

不过，第7节研究的代数理论的一个重要特点在 \SGARefLink{sga7.sga7i.eix.n13-4-17}{13.4.17} 的描述中缺失。假设 $A$ 在 $S$ 的一个无处稠密解析子集 $Z$ 的补开集 $U$ 上为类阿贝尔群——例如 $\dim S=1$ 时 $Z=\{s\}$——于是由（\SGARefLink{sga7.sga7i.eix.n13-1-5}{13.1.5}），$R|U$ 是 $U$ 上的局部系统。那么 $U$ 上的局部系统 $R/R^t$ 应当是\underline{常值的}；更确切地说，它应与局部系统 ${R'}^f$ 对偶，其中 $R'$ 是某个适当的 $S$ 上解析群 $A'$ 所对应的、在 $S$ 上平展的解析群（如同 $R$ 对应于 $A$），使 $A'|U$ 为 $A|U$ 的对偶类阿贝尔群。用第5节的记号，其代数对应物是（\SGARefLink{sga7.sga7i.eix.n5-5-9}{5.5.9}）中 $T_\ell(A_K)/T_\ell(A_K)^t$ 与 $T_\ell(A'_K)^f$ 之间的对偶性，它来自对偶性定理 \SGARefLink{sga7.sga7i.eix.n5-2}{5.2}。

我不知道仅在 \SGARefLink{sga7.sga7i.eix.n13-4-17}{13.4.17} 的假设下，这种断言是否成立，甚至不知道对 $t\in U$，$\pi_1(U,t)$ 在 $R_t$ 上的作用是否幺幂。本段使用的朴素方法似乎也不能证明这样的结果，即使另加 $A|U$ 带有“极化”的假设，即带有在每个纤维上丰富的可逆层。然而，“半单群方面的人”似乎知道如何证明它，至少当 $S$ 为非奇异的 $1$ 维空间时如此；参见报告 [9]。下一小节将说明，在 $A$ 的相对代数性假设下，一切均可成立，同时把此结果与第7节的构造联系起来。

% END INCLUDED SGA7I_English_current_buildable_source_20260823/components/180_expose_IX_section_13_4.tex

% BEGIN INCLUDED SGA7I_English_current_buildable_source_20260823/components/181_expose_IX_section_13_5.tex
\medskip
\noindent
\SGARefTarget{sga7.sga7i.eix.n13-5}{13.5}. \underline{相对代数情形。}
现在假设 $A$ 来自 $S$ 上的一个光滑分离相对群概形，其纤维为阿贝尔概形对环面的扩张。我们以同一字母 $A$ 表示此相对概形。（允许自己这样滥用记号，是因为从所考虑类型的相对群概形范畴到 $S$ 上解析群范畴的函子 $G\longmapsto G^{\mathrm{an}}$ 是全忠实的，用相对 GAGA 型结果 [10] 容易证明这一点。）关于相对概形的概念及 $S$ 上相对概形与 $S$ 上解析空间之间的关系，均参见 [10]（也可参见 [2，第15讲]）。

若 $\underline O$ 表示 $S$ 在 $s$ 处的局部环，则相对概形 $A$ 在 $s$ 附近由普通概形 $A_{\underline O}$ 确定，后者是 $A$ 在 $s$ 处局部化得到的 $\operatorname{Spec}(\underline O)$ 上的概形。于是显然，$A_{\underline O}$ 沿其特殊纤维（同构于 $A_s=A_o$）的形式完备化，同构于解析群 $A$ 沿 $A_s$ 的形式完备化（\SGARefLink{sga7.sga7i.eix.n13-2}{13.2}）；这里按前面引入的记号约定，把 $A_n=A\times_S S_n$ 等同于 $S$ 上的代数概形。在 $\hat A$ 作为 $\hat{\underline O}$ 上的形式群概形可代数化（\SGARefLink{sga7.sga7i.eix.n7-2-1}{7.2.1}）的有利情形下——特别地，当 $S$ 在 $s$ 处正规而 $A$ 在 $S$ 上拟射影时便如此（\SGARefLink{sga7.sga7i.eix.n7-1-5}{7.1.5}）——于是得到 $\operatorname{Spec}(\hat{\underline O})$ 上的群概形 $A^{\natural}_{\hat{\underline O}}$，它是阿贝尔概形 $B_{\hat{\underline O}}$ 对环面的扩张。为简单起见，假设 $B_{\hat{\underline O}}$ 在 $\hat{\underline O}$ 上射影（我们知道，当 $\underline O$，从而 $\hat{\underline O}$ 正规时的确如此）。我们将证明：

\medskip
\noindent
\underline{命题 \SGARefTarget{sga7.sga7i.eix.n13-5-1}{13.5.1}.} \underline{在上述条件下，（\SGARefLink{sga7.sga7i.eix.n13-3-3}{13.3.3}）中的解析群 $A^{\natural}$ 在 $s$ 的一个邻域中，来自一个在唯一同构意义下确定的相对群概形，它是 $S$ 上一个射影相对阿贝尔概形对环面的扩张。若 $A^{\natural}_{\underline O}$ 表示它在 $\underline O=\underline O_{S,s}$ 上定义的群概形，则群概形}

\begin{equation*}
 A^{\natural}_{\underline O}\otimes_{\underline O}
 \hat{\underline O}
\end{equation*}

\noindent
\underline{连同其扩张结构，典范同构于上述 Raynaud 扩张 $A^{\natural}_{\hat{\underline O}}$。}

后一个断言由前一个断言，结合 \SGARefLink{sga7.sga7i.eix.n7-2-1}{7.2.1} 立即得到，因为由（\SGARefLink{sga7.sga7i.eix.n13-3-6}{13.3.6}），$A^{\natural}_{\underline O}\otimes_{\underline O}\hat{\underline O}$ 沿其特殊纤维的形式完备化为 $\hat A_{\underline O}$。为证明 $A^{\natural}$ 在 $s$ 附近的代数性，按 \SGARefLink{sga7.sga7i.eix.n7-2-1}{7.2.1} 的方法，立即归结为证明 $A^{\natural}$ 的阿贝尔商 $B$ 的代数性，而这来自下一个引理。

\medskip
\noindent
\underline{引理 \SGARefTarget{sga7.sga7i.eix.n13-5-2}{13.5.2}.} \underline{设 $B$ 是解析空间 $S$ 上的类阿贝尔解析群，$s$ 是 $S$ 的一点，使沿纤维 $B_s$ 的完备化 $\hat B$ 可代数化；即 $B_n=B\times_S S_n$ 可代数化，而且将 $\hat B$ 视为 $\hat{\underline O}_{S,s}$ 上的形式概形时它可代数化，或等价地，来自 $\hat{\underline O}$ 上的阿贝尔概形 $B_{\hat{\underline O}}$。再假设 $B_{\hat{\underline O}}$ 在 $\hat{\underline O}$ 上射影。则 $B$ 在 $s$ 的一个邻域中在 $S$ 上射影。}

假设意味着存在形式解析空间 $\hat B$ 的极化，即 $B_n$ 的一族可以“粘合”的极化 $p_n$。结论意味着在 $s$ 的一个邻域上存在 $B$ 的极化 $p$。设 $P$ 为 $B$ 在 $S$ 上的相对极化所成的 $S$ 上解析空间，它表示 $S$ 上解析空间 $S'$ 的范畴上的函子

\begin{equation*}
 S'\longmapsto
 \{\text{$B_{S'}$ 在 $S'$ 上的极化}\}
\end{equation*}

\noindent
这是 $S$ 上 N\'eron--Severi 函子 $\underline{\mathrm{NS}}_{B/S}$ 的开子函子。最简单的标准论证（在此省略）说明，$\underline{\mathrm{NS}}_{B/S}$，从而 $P$，由 $S$ 上的一个\underline{非分歧}解析空间表示。于是，假设是说存在 $P$ 在 $S$ 上于 $s$ 处的一个\underline{形式}截面 $p$。由于 $P$ 在 $S$ 上非分歧，故局部拟有限，由 [2，第18讲，第1节] 可知，$P$ 在 $S$ 上于 $s$ 处的每个形式截面都来自 $s$ 附近的一个解析截面，即 $B$ 在 $s$ 的一个邻域上的一个极化。断言得证。

\medskip
\noindent
\underline{注 \SGARefTarget{sga7.sga7i.eix.n13-5-3}{13.5.3}.}

\smallskip
\noindent
a) 若在 \SGARefLink{sga7.sga7i.eix.n13-5-2}{13.5.2} 中删去 $B_{\hat{\underline O}}$ 在 $\hat{\underline O}$ 上射影的假设，\footnote{印本此处为“在 $S$ 上射影”；引理 \SGARefLink{sga7.sga7i.eix.n13-5-2}{13.5.2} 的假设和随后的正规性论证要求“在 $\hat{\underline O}$ 上射影”。}仍可得到 $B$ 在 $s$ 附近可代数化。为此使用一个未发表的结果（M.~Artin 和 P.~Deligne 分别独立得到，使用 VI 2 的层 $\underline T^i$ 的理论）：若 $X$ 在 $S$ 上紧合且 $X_{\mathrm{red}}$ 可代数化，则 $X$ 可代数化。由此归结为 $S$ 约化的情形。另一方面，当 $S$ 正规，从而 $\hat{\underline O}$ 正规时，已知 $\hat{\underline O}$ 上的每个阿贝尔概形都在 $\hat{\underline O}$ 上射影（例如参见 [23]），故可应用13.5.2。为得到任意约化 $S$ 的情形，可从 $S$ 的正规化 $S'$ 作非平坦下降，使用如下（微妙的）结果，在此不证明：诺特概形的紧合满态射 $f:S'\longrightarrow S$（即使 $\underline O_S\longrightarrow f_*(\underline O_{S'})$ 单射的态射），是相对阿贝尔概形所成纤维范畴的一个\underline{有效下降态射}。

\smallskip
\noindent
b) 看来，\SGARefLink{sga7.sga7i.eix.n13-5-2}{13.5.2} 的证明也应当能证明 $S$ 上每个紧合解析空间 $B$ 的类似结果。可能更微妙的问题是：若 $S$ 上紧合的解析空间 $X$ 沿 $X_s$ 的形式完备化 $\hat X$ 可代数化，则其本身在 $s$ 附近是否在 $S$ 上可代数化。只指出，此处 $X$ 在 $S$ 上的相对可代数化性的含义应比 [10] 中的更广：应要求 $X$ 在 $s$ 附近来自 $\underline O_{S,s}$ 上一个 M.~Artin [1] 意义下的\underline{代数空间}。

\medskip
\noindent
\SGARefTarget{sga7.sga7i.eix.n13-5-4}{13.5.4}. 置身于 \SGARefLink{sga7.sga7i.eix.n13-5-1}{13.5.1} 的条件下。对每个整数 $n>0$，还选取 $\mathbb C$ 上的标准群同构

\begin{equation*}
 \mu_{n\mathbb C}\simeq(\mathbb Z/n\mathbb Z)_{\mathbb C},
 \qquad m\longmapsto\exp(2\pi i m/n)\qquad .
\end{equation*}

\noindent
对固定素数 $\ell$ 的各次幂取投射极限，得到 $\mathbb C$ 上 $\ell$ 进层的典范同构

\begin{equation*}
 \mathbb Z_\ell(1)_{\mathbb C}\simeq\mathbb Z_{\ell\mathbb C}
 \qquad ,
\end{equation*}

\noindent
我们将利用它等同两端。

将 $\underline{\operatorname{Ext}}^i_{\mathbb Z}(\mathbb Z/n\mathbb Z,-)$ 的正合列应用于正合列（\SGARefLink{sga7.sga7i.eix.n13-1-3}{13.1.3}），得到典范同构

\begin{equation}\tag{\SGARefTarget{sga7.sga7i.eix.n13-5-5}{13.5.5}}
 {}_nA^{\mathrm o}\simeq
 \underline R\otimes_{\mathbb Z}\mathbb Z/n\mathbb Z
 \qquad ,
\end{equation}

\noindent
当 $n$ 取遍给定素数 $\ell$ 的各次幂时，它给出 $S$ 上群的投射系统的同构

\begin{equation}\tag{\SGARefTarget{sga7.sga7i.eix.n13-5-6}{13.5.6}}
 T_\ell(A^{\mathrm o})\simeq
 \underline R\otimes_{\mathbb Z}\mathbb Z_\ell
 \qquad ,
\end{equation}

\noindent
其中把 $\mathbb Z_\ell$ 看作 $\mathbb Z/\ell^{n+1}\mathbb Z$ 所成的投射系统。在这里考虑的情形中，$A$ 来自 $S$ 上的分离群概形，于是立即可见 $\underline R\otimes\mathbb Z/n\mathbb Z$ 可看作相对概形。这使符号 $(\underline R\otimes_{\mathbb Z}\mathbb Z_\ell)_{\underline O}$ 有了意义（而 $R_{\underline O}$ 一般没有意义！），并得到典范同构

\begin{equation}\tag{\SGARefTarget{sga7.sga7i.eix.n13-5-7}{13.5.7}}
 T_\ell(A_{\underline O})\simeq
 (\underline R\otimes_{\mathbb Z}\mathbb Z_\ell)_{\underline O}
 \qquad .
\end{equation}

\noindent
由第2段和 \SGARefLink{sga7.sga7i.eix.n13-3}{13.3} 的构造立即得到，上述同构将不变部分映为不变部分，将环面部分映为环面部分：

\begin{equation}\tag{\SGARefTarget{sga7.sga7i.eix.n13-5-8}{13.5.8}}
 T_\ell(A_{\underline O})^f\simeq
 (\underline R^f\otimes_{\mathbb Z}\mathbb Z_\ell)_{\underline O},
 \qquad
 T_\ell(A_{\underline O})^t\simeq
 (\underline R^t\otimes_{\mathbb Z}\mathbb Z_\ell)_{\underline O}
 \qquad .
\end{equation}

\noindent
若 $A'$ 是 $S$ 上的相对解析群，满足与 $A$ 相同的假设，而 $W_{\underline O}$ 是 $(A_{\underline O},A'_{\underline O})$ 对 $\mathbb G_{m\underline O}$ 的双扩张，因而在 $s$ 的一个邻域上定义 $(A,A')$ 对 $\mathbb G_{mS}$ 的双扩张（现在使用复解析意义的概念，即在赋予适当拓扑的 $S$ 上解析空间景中工作而得到的概念），则由 VIII \SGARefLink{sga7.sga7i.eviii.n2-2}{2.2}，$W$ 定义配对

\begin{equation}\tag{\SGARefTarget{sga7.sga7i.eix.n13-5-9}{13.5.9}}
 (\underline R\otimes\mathbb Z_\ell)\times
 (\underline R'\otimes_{\mathbb Z}\mathbb Z_\ell)
 \simeq T_\ell(A)\times T_\ell(A')
 \longrightarrow\mathbb Z_\ell(1)_S
 \simeq\mathbb Z_{\ell S}\qquad ,
\end{equation}

\noindent
同样显然，相应的配对

\begin{equation}\tag{\SGARefTarget{sga7.sga7i.eix.n13-5-10}{13.5.10}}
 T_\ell(A_{\underline O})\times T_\ell(A'_{\underline O})
 \longrightarrow\mathbb Z_\ell(1)_S
\end{equation}

\noindent
正是所引文献联系于 $(A_{\underline O},A'_{\underline O})$ 的双扩张 $W_{\underline O}$ 的配对。也不难看出，限制到 $s$ 的一个邻域后，配对（\SGARefLink{sga7.sga7i.eix.n13-5-9}{13.5.9}）来自唯一的一个（不再依赖于 $\ell$ 的）配对

\begin{equation}\tag{\SGARefTarget{sga7.sga7i.eix.n13-5-11}{13.5.11}}
 \underline R\otimes\underline R'\longrightarrow\mathbb Z
 \qquad .
\end{equation}

\noindent
有了这些准备，有关 $T_\ell(A_{\underline O})$ 的代数结果可用于推出关于层 $\underline R$ 的结论。因此，配对（\SGARefLink{sga7.sga7i.eix.n13-5-11}{13.5.11}）在 $\underline R^t\times\underline R'{}^f$ 及 $\underline R^f\times\underline R'{}^t$ 上为零（\SGARefLink{sga7.sga7i.eix.n7-4-8}{7.4.8}）。此外，若 $S$ 是非奇异的 $1$ 维空间，而在 $U=S-\{s\}$ 上，$A$ 和 $A'$ 是通过 $W$ 彼此对偶的相对阿贝尔概形，则在 $s$ 的一个充分小的邻域内，$\underline R^t_U$ 和 $\underline R'{}^f_U$，以及 $\underline R^f_U$ 和 $\underline R'{}^t_U$，分别互为精确零化子（\SGARefLink{sga7.sga7i.eix.n7-4-8}{7.4.8}）。由于（\SGARefLink{sga7.sga7i.eix.n13-5-11}{13.5.11}）在 $U$ 上为完美对偶，故在 $U$ 上得到 $\underline R/\underline R^t$ 与 $\underline R'{}^f$ 之间的对偶性。特别地，$\underline R/\underline R^t$ 为常值的，因而 $\pi_1(U,t)$ 在 $\underline R_t$ 上的作用更是 $2$ 级幺幂的。类似地，在 $U$ 上有 $\underline R/\underline R^f$ 与 $\underline R'{}^t$ 之间的完美对偶。由于后者还同构于 $S$ 上格 $M_S$ 的对偶环面，其中 $M$ 是（\SGARefLink{sga7.sga7i.eix.n5-5-4}{5.5.4}）中定义的群，得到典范同构

\begin{equation}\tag{\SGARefTarget{sga7.sga7i.eix.n13-5-12}{13.5.12}}
 \underline R_U/\underline R^f_U\simeq M_U
 \qquad ,
\end{equation}

\noindent
而（\SGARefLink{sga7.sga7i.eix.n13-3-4}{13.3.4}）中的同态 $i$ 可写为同态

\begin{equation}\tag{\SGARefTarget{sga7.sga7i.eix.n13-5-13}{13.5.13}}
 i:M_U\longrightarrow A^\natural
 \qquad ,
\end{equation}

\noindent
下一段将在稍有不同的背景下再次遇到它。

% END INCLUDED SGA7I_English_current_buildable_source_20260823/components/181_expose_IX_section_13_5.tex


% BEGIN INCLUDED SGA7I_English_current_buildable_source_20260823/components/182_expose_IX_section_14_1.tex
\medskip
\noindent
14. \underline{与超越理论的联系：刚性解析情形。典范同态}
$\underline M_K\longrightarrow A_K^{\natural\mathrm o}$.

\medskip
\noindent
\SGARefTarget{sga7.sga7i.eix.n14-1}{14.1}. 此处假定已知刚性解析空间理论的基础，可参阅 [12]、[13] 和 [34]。

设 $S$ 是完备性状，$A_K$ 是其分式域上的阿贝尔概形。在 \SGARefLink{sga7.sga7i.eix.n8-1}{8.1} 中，从半稳定情形的 Raynaud 扩张作下降，构造了一个扩张 $G_K=A_K^{\natural\mathrm o}$（\SGARefLink{sga7.sga7i.eix.n8-1-1}{8.1.1}）

\begin{equation}\tag{\SGARefTarget{sga7.sga7i.eix.n14-1-1}{14.1.1}}
 0\longrightarrow T_K\longrightarrow G_K\longrightarrow B_K
 \longrightarrow0
\end{equation}

\noindent
即阿贝尔概形 $B_K$ 对环面 $T_K$ 的扩张。考虑对偶阿贝尔概形 $A'_K$，同样有扩张 $G'_K=A'{}^{\natural\mathrm o}_K$

\begin{equation}\tag{\SGARefTarget{sga7.sga7i.eix.n14-1-2}{14.1.2}}
 0\longrightarrow T'_K\longrightarrow G'_K\longrightarrow B'_K
 \longrightarrow0\qquad ,
\end{equation}

\noindent
其中还把 $B'_K$ 等同于 $B_K$ 的对偶阿贝尔概形（\SGARefLink{sga7.sga7i.eix.n8-2}{8.2}）。如（\SGARefLink{sga7.sga7i.eix.n8-1-8}{8.1.8}），以 $\underline M_K$ 和 $\underline M'_K$ 分别表示 $K$ 上与 $T'_K$ 和 $T_K$ 对偶的格：

\begin{equation}\tag{\SGARefTarget{sga7.sga7i.eix.n14-1-3}{14.1.3}}
 \underline M_K=D(T'_K),\qquad \underline M'_K=D(T_K)\qquad .
\end{equation}

\noindent
于是扩张结构（\SGARefLink{sga7.sga7i.eix.n14-1-2}{14.1.2}）和（\SGARefLink{sga7.sga7i.eix.n14-1-1}{14.1.1}）分别可解释为（VIII \SGARefLink{sga7.sga7i.eviii.n3-7}{3.7}）等价于同态的数据

\begin{equation}\tag{\SGARefTarget{sga7.sga7i.eix.n14-1-4}{14.1.4}}
 \underline M_K\xrightarrow{\ p_K\ }B_K,\qquad
 \underline M'_K\xrightarrow{\ p'_K\ }B'_K\qquad .
\end{equation}

设 $W_K^{\natural}$ 是定义这些阿贝尔概形之间典范对偶性（\SGARefLink{sga7.sga7i.eix.n8-2}{8.2}）的 $(B_K,B'_K)$ 对 $\mathbb G_{mK}$ 的双扩张，$E_K$ 是把 $W_K^{\natural}$ 沿同态（\SGARefLink{sga7.sga7i.eix.n14-1-4}{14.1.4}）拉回得到的 $(\underline M_K,\underline M'_K)$ 对 $\mathbb G_{mK}$ 的双扩张。立即可验证，$E_K$ 的平凡化典范地一一对应于同态

\begin{equation}\tag{\SGARefTarget{sga7.sga7i.eix.n14-1-5}{14.1.5}}
 i:\underline M_K\longrightarrow G_K
\end{equation}

\noindent
它们提升 $p_K:\underline M_K\longrightarrow B_K$；或对称地，对应于同态

\begin{equation}\tag{\SGARefTarget{sga7.sga7i.eix.n14-1-6}{14.1.6}}
 i':\underline M'_K\longrightarrow G'_K
\end{equation}

\noindent
它们提升 $p'_K:\underline M'_K\longrightarrow B'_K$。

在此理解下，用上标 ${}^{\mathrm{an}}$ 表示所考虑代数群所联系的 $K$ 上的刚性解析群 [13]。[24] 中定义一个典范正合列

\begin{equation}\tag{\SGARefTarget{sga7.sga7i.eix.n14-1-7}{14.1.7}}
 0\longrightarrow\underline M_K^{\mathrm{an}}
 \xrightarrow{\ i^{\mathrm{an}}\ }G_K^{\mathrm{an}}
 \xrightarrow{\ e\ }A_K^{\mathrm{an}}\longrightarrow0\qquad .
\end{equation}

\noindent
这里 $i^{\mathrm{an}}$ 也可解释为通常群概形的同态 $i:\underline M_K\longrightarrow G_K$，它\underline{提升}（\SGARefLink{sga7.sga7i.eix.n14-1-4}{14.1.4}）中的\underline{同态} $p$，即它具有（\SGARefLink{sga7.sga7i.eix.n14-1-5}{14.1.5}）考虑的形式。这个正合列是正合列（\SGARefLink{sga7.sga7i.eix.n13-3-4}{13.3.4}）的刚性解析对应物。它对 $A_K$ 具有函子性，其形成与每个性状变换 $S'\longrightarrow S$ 相容。

后一个事实还说明，为定义和研究（\SGARefLink{sga7.sga7i.eix.n14-1-7}{14.1.7}），利用有限下降，可归结到 $A_K$ 有半稳定约化的情形。此时，群概形 $T_K,G_K,B_K,\underline M_K,\ldots$ 是 $S$ 上群概形 $T,G,B,M,\ldots$ 的泛纤维，其中 $G$ 是阿贝尔概形 $B$ 对环面 $T$ 的扩张（第7段）。于是可考虑 $G$ 沿其特殊纤维的形式完备化 $\hat G$，由构造有典范同构（\SGARefLink{sga7.sga7i.eix.n7-2-6}{7.2.6}）

\begin{equation}\tag{\SGARefTarget{sga7.sga7i.eix.n14-1-8}{14.1.8}}
 \hat G\xrightarrow{\ \sim\ }\widehat{A^{\mathrm o}}
 \qquad .
\end{equation}

\noindent
转到“泛纤维”，它们是 $K$ 上的刚性解析群，得到典范同构

\begin{equation}\tag{\SGARefTarget{sga7.sga7i.eix.n14-1-9}{14.1.9}}
 \hat G^{\mathrm{an}}\xrightarrow{\ \sim\ }
 \widehat{A^{\mathrm o}}\,{}^{\mathrm{an}}
 \qquad .
\end{equation}

\noindent
另一方面，有典范同态（刚性解析开浸入，但一般不是同构）

\begin{equation}\tag{\SGARefTarget{sga7.sga7i.eix.n14-1-10}{14.1.10}}
 \hat G^{\mathrm{an}}\lhook\joinrel\longrightarrow G_K^{\mathrm{an}},
 \qquad
 \widehat{A^{\mathrm o}}\,{}^{\mathrm{an}}
 \lhook\joinrel\longrightarrow A_K^{\mathrm{an}}
 \qquad .
\end{equation}

\noindent
（\SGARefLink{sga7.sga7i.eix.n14-1-7}{14.1.7}）中的同态 $e$ 由图表的交换性唯一刻画：

\begin{equation}\tag{\SGARefTarget{sga7.sga7i.eix.n14-1-11}{14.1.11}}
\begin{tikzcd}
 \hat G^{\mathrm{an}} \arrow[r] \arrow[d]
   & \widehat{A^{\mathrm o}}\,{}^{\mathrm{an}} \arrow[d] \\
 G_K^{\mathrm{an}} \arrow[r] & A_K^{\mathrm{an}}
\end{tikzcd}
\end{equation}

\noindent
其中未标记的箭头就是（\SGARefLink{sga7.sga7i.eix.n14-1-9}{14.1.9}）和（\SGARefLink{sga7.sga7i.eix.n14-1-10}{14.1.10}）中的箭头。

\medskip
\noindent
\underline{注 \SGARefTarget{sga7.sga7i.eix.n14-1-12}{14.1.12}.}

\smallskip
\noindent
a) 注意，使上述图表交换的 $e$ 的唯一性是直接的，因为刚性解析空间 $G_K^{\mathrm{an}}$ 连通。另一方面，$e$ 的存在性远非平凡，在此仅请读者参阅所引 M.~Raynaud 的札记。此外，所引文献中给出的 $i^{\mathrm{an}}$ 的定义同样具有超越性质，尽管在通过第7段的\underline{形式几何}构造得到 $G$ 本身，从而得到 $G_K$ 之后，可将它解释为一个本质上为代数性的对象（\SGARefLink{sga7.sga7i.eix.n14-1-5}{14.1.5}）。我不知道它的纯代数几何描述，但下面（\SGARefLink{sga7.sga7i.eix.n14-4}{14.4}）将提出一个“算术”刻画，通过归结到环 $\mathbb Z$ 上有限型的情形得到。因此，从代数几何的角度，同态（\SGARefLink{sga7.sga7i.eix.n14-1-5}{14.1.5}）似乎是比前几段遇到的结构更隐蔽的一项结构。

\smallskip
\noindent
b) 如前所述，同态 $i$ 典范地联系于 $(\underline M_K,\underline M'_K)$ 的双扩张 $E_K$ 的一个确定的平凡化，或等价地，联系于提升 $p'_K$ 的一个确定的同态（\SGARefLink{sga7.sga7i.eix.n14-1-6}{14.1.6}）。于是可证，后者恰好是正合列

\begin{equation}\tag{\SGARefTarget{sga7.sga7i.eix.n14-1-13}{14.1.13}}
 0\longrightarrow\underline M'{}^{\mathrm{an}}_K
 \xrightarrow{\ i'{}^{\mathrm{an}}\ }G'{}^{\mathrm{an}}_K
 \longrightarrow A'{}^{\mathrm{an}}_K\longrightarrow0
\end{equation}

\noindent
中出现的同态；此列类似于（\SGARefLink{sga7.sga7i.eix.n14-1-7}{14.1.7}），联系于 $A_K$ 的对偶阿贝尔概形 $A'_K$（除非应取其相反同态？？？）。

% END INCLUDED SGA7I_English_current_buildable_source_20260823/components/182_expose_IX_section_14_1.tex

% BEGIN INCLUDED SGA7I_English_current_buildable_source_20260823/components/183_expose_IX_section_14_2.tex
\par\medskip
\noindent
\SGARefTarget{sga7.sga7i.eix.n14-2}{14.2}. \underline{同态}
$i:\underline M_K\longrightarrow G_K$ \underline{与单值作用的关系。}

设 $n$ 是整数 $>0$。由正合列（\SGARefLink{sga7.sga7i.eix.n14-1-7}{14.1.7}），利用 $K$ 上刚性解析空间景中的 $\underline{\operatorname{Ext}}^i(\mathbb Z/n\mathbb Z,-)$ 的正合列，得到典范正合列（考虑到 $\underline M_K^{\mathrm{an}}$ 无挠而 $G_K^{\mathrm{an}}$ 是 $n$-可除的）：

\begin{equation*}
 0\longrightarrow{}_nG_K^{\mathrm{an}}
 \longrightarrow{}_nA_K^{\mathrm{an}}
 \longrightarrow\underline M_{K,n}^{\mathrm{an}}
 \longrightarrow0\qquad ,
\end{equation*}

\noindent
其中 $\underline M_{K,n}\overset{\mathrm{dfn}}{=}\underline M_K/n\underline M_K$。由于此正合列中出现的各项都来自 $K$ 上的有限群概形，可将它重写为来自一个正合列

\begin{equation}\tag{\SGARefTarget{sga7.sga7i.eix.n14-2-1}{14.2.1}}
 0\longrightarrow{}_nG_K\longrightarrow{}_nA_K
 \longrightarrow\underline M_{K,n}\longrightarrow0\qquad .
\end{equation}

\noindent
当 $n$ 变化时，这些正合列具有预期的行为；对素数 $\ell$ 的各次幂取投射极限，得到类似的正合列

\begin{equation}\tag{\SGARefTarget{sga7.sga7i.eix.n14-2-2}{14.2.2}}
 0\longrightarrow T_\ell(G_K)\longrightarrow T_\ell(A_K)
 \longrightarrow\underline M_K\otimes_{\mathbb Z}\mathbb Z_\ell
 \longrightarrow0\qquad .
\end{equation}

\noindent
但纯代数方法已给出这样的正合列：借助第一个公式（\SGARefLink{sga7.sga7i.eix.n8-1-2}{8.1.2}）及（\SGARefLink{sga7.sga7i.eix.n8-1-7}{8.1.7}），有

\begin{equation}\tag{\SGARefTarget{sga7.sga7i.eix.n14-2-3}{14.2.3}}
 T_\ell(G_K)\xrightarrow{\ \sim\ }T_\ell(A_K)^{\mathrm{ef}}
 \subset T_\ell(A_K),\qquad
 T_\ell(A_K)/T_\ell(A_K)^{\mathrm{ef}}
 \simeq\underline M_K\otimes_{\mathbb Z}\mathbb Z_\ell\qquad .
\end{equation}

\noindent
在此理解下，我们断言，\underline{代数理论中得到的同态}（\SGARefLink{sga7.sga7i.eix.n14-2-3}{14.2.3}）\underline{恰好就是由超越理论中的正合列}（\SGARefLink{sga7.sga7i.eix.n14-2-2}{14.2.2}）\underline{得到的同态}。对于第一个同态，即对应于（\SGARefLink{sga7.sga7i.eix.n14-2-2}{14.2.2}）中第一个同态的同态，归结到半稳定情形后，我们的断言立即来自（\SGARefLink{sga7.sga7i.eix.n14-1-11}{14.1.11}）的交换性，因为那里的竖直箭头在 $T_\ell$ 上诱导同构，而将 $T_\ell$ 应用于水平箭头，分别得到（\SGARefLink{sga7.sga7i.eix.n14-2-3}{14.2.3}）和（\SGARefLink{sga7.sga7i.eix.n14-2-2}{14.2.2}）中所讨论的箭头。对 $T_\ell(A_K)/T_\ell(A_K)^{\mathrm{ef}}$ 与 $\underline M_K\otimes_{\mathbb Z}\mathbb Z_\ell$ 之间所得到的两个同构，我们承认所述相容性。此外注意，\underline{这一相容性}（对给定的 $\ell$）\underline{已足以用超越方法唯一刻画}（\SGARefLink{sga7.sga7i.eix.n14-1-7}{14.1.7}）\underline{中出现的同态}（\SGARefLink{sga7.sga7i.eix.n14-1-5}{14.1.5}），将其解释为 $\underline M_K^{\mathrm{an}}$ 到 $\operatorname{Ker}e$ 的同构。

也可说，（\SGARefLink{sga7.sga7i.eix.n14-1-7}{14.1.7}）中的同态 $e$ 将各个 $T_\ell(A_K)/T_\ell(A_K)^{\mathrm{ef}}$ 等同于 $K$ 上同一个格 $\underline R_K=\operatorname{Ker}e$ 的各个 $\ell$ 进完备化 $\underline R_K\otimes\mathbb Z_\ell$，而前述 pro-$\ell$ 群系统上的这个“整结构”，在显然意义下，与公式（\SGARefLink{sga7.sga7i.eix.n8-1-7}{8.1.7}）（正交性定理 \SGARefLink{sga7.sga7i.eix.n5-2}{5.2}）给出的整结构相容。于是典范地定义同构

\begin{equation*}
 i:\underline M_K\xrightarrow{\ \sim\ }\operatorname{Ker}e
 \qquad .
\end{equation*}

由于同态 $i$（\SGARefLink{sga7.sga7i.eix.n14-1-5}{14.1.5}）使我们能够重建 $T_\ell(A_K)$ 上通常的扩张结构，即（\SGARefLink{sga7.sga7i.eix.n14-2-2}{14.2.2}），从而借助由（\SGARefLink{sga7.sga7i.eix.n14-1-1}{14.1.1}）得到的正合列也能重建 \SGARefLink{sga7.sga7i.eix.n9-6}{9.6} 中描述的\underline{混合扩张}结构，

\begin{equation*}
\begin{array}{ccccccccc}
0&\longrightarrow&T_\ell(T_K)&\longrightarrow&T_\ell(G_K)
 &\longrightarrow&T_\ell(B_K)&\longrightarrow&0\\[5ex]
 &&\underline M'_K\otimes\mathbb Z_\ell(1)&&&&&&
\end{array}
\end{equation*}

\noindent
所以在半稳定情形下，它还能重建\underline{单值配对}（\SGARefLink{sga7.sga7i.eix.n9-1-2}{9.1.2}）

\begin{equation}\tag{\SGARefTarget{sga7.sga7i.eix.n14-2-4}{14.2.4}}
 u_\ell:\underline M\otimes\underline M'
 \longrightarrow\mathbb Z_\ell\qquad ,
\end{equation}

\noindent
这些配对在 \SGARefLink{sga7.sga7i.eix.n9-6}{9.6} 中正是用混合扩张定义的。这里只指出将 \SGARefLink{sga7.sga7i.eix.n9-6}{9.6} 的构造改写后的结果。为此注意，在半稳定情形下，$p_K$ 和 $p'_K$ 来自同态

\begin{equation*}
 p:\underline M\longrightarrow B,\qquad
 p':\underline M'\longrightarrow B'\qquad ,
\end{equation*}

\noindent
其中 $B,B'$ 是 $S$ 上彼此对偶的阿贝尔概形。于是沿 $(p,p')$ 拉回，得到 $(\underline M,\underline M')$ 对 $\mathbb G_{mS}$ 的双扩张 $E$，其泛纤维为 \SGARefLink{sga7.sga7i.eix.n4-1}{4.1} 中已考虑过的双扩张 $E_K$。

为简单起见，假设 $\underline M$ 和 $\underline M'$ 为常值的，取值分别为 $M$ 和 $M'$。对每对 $(m,m')\in M\times M'$，$E_{m,m'}$ 是 $\mathbb G_{mS}$-挠子，故对应于 $S$ 的环 $V$ 上的可逆模 $L_{m,m'}$。所以它有一个基 $u_0$，在乘以 $V$ 的单位的意义下确定。因此得到典范的“赋值”映射

\begin{equation*}
 \underline v:E_{m,m'}(K)\longrightarrow\mathbb Z
\end{equation*}

\noindent
它把 $L_{m,m'}\otimes_V K$ 的每个基 $u$，写为 $u=\lambda u_0$（$\lambda\in K^*$），映为元素

\begin{equation*}
 \underline v(u)=\underline v(\lambda)\in\mathbb Z\qquad ,
\end{equation*}

\noindent
其中右端的 $\underline v$ 表示与 $V$ 联系的赋值（按通常方式正规化，使值群为 $\mathbb Z$）；所得整数显然与 $u_0$ 的选取无关。若 $\varphi$ 现在是 $E$ 的一个平凡化，可考虑相应映射

\begin{equation}\tag{\SGARefTarget{sga7.sga7i.eix.n14-2-5}{14.2.5}}
 (m,m')\longmapsto\underline v\bigl(\varphi(m,m')\bigr):
 M\times M'\longrightarrow\mathbb Z\qquad .
\end{equation}

\noindent
在此理解下，发现\underline{单值同态}（\SGARefLink{sga7.sga7i.eix.n14-2-4}{14.2.4}）\underline{取值于 $\mathbb Z$，且恰好为同态}（\SGARefLink{sga7.sga7i.eix.n14-2-5}{14.2.5}），\underline{其中 $\varphi$ 是与同态 $i$}（\SGARefLink{sga7.sga7i.eix.n14-1-5}{14.1.5}）\underline{联系的 $E_K$ 的分裂}。

\medskip
\noindent
\underline{注 \SGARefTarget{sga7.sga7i.eix.n14-2-6}{14.2.6}.} 特别地，（借助超越的刚性解析理论）重新得到单值配对取整值且与 $\ell$ 无关的事实（\SGARefLink{sga7.sga7i.eix.n10-4}{10.4} a)），它已在 \SGARefLink{sga7.sga7i.eix.n12-5}{12.5} 中用另一种（同样超越的）方法证明。据 M.~Raynaud 所言，刚性解析理论也重新得到正性断言 \SGARefLink{sga7.sga7i.eix.n10-4}{10.4} b)（这也包含在 \SGARefLink{sga7.sga7i.eix.n12-5}{12.5} 中）。

\medskip
\noindent
\underline{注 \SGARefTarget{sga7.sga7i.eix.n14-2-7}{14.2.7}.} 将 $\underline M_K\otimes_{\mathbb Z}\mathbb Q$ 看作 $\underline M_K\otimes_{\mathbb Z}\mathbb Q/\mathbb Z$ 对 $\underline M_K$ 的扩张。利用同态（\SGARefLink{sga7.sga7i.eix.n14-1-5}{14.1.5}）$i:\underline M_K\longrightarrow G_K$，得到一个扩张 $A_K^b$

\begin{equation}\tag{\SGARefTarget{sga7.sga7i.eix.n14-2-8}{14.2.8}}
 0\longrightarrow G_K\longrightarrow A_K^b
 \longrightarrow\underline M_K\otimes\mathbb Q/\mathbb Z
 \longrightarrow0\qquad .
\end{equation}

\noindent
还须注意，由于 $G_K$ 是可除群，给出这样的扩张，等价于给出一族（\SGARefLink{sga7.sga7i.eix.n14-2-2}{14.2.2}）形式的扩张，即对（\SGARefLink{sga7.sga7i.eix.n14-2-8}{14.2.8}）应用各个 $T_\ell$ 所得的扩张。由于 $i$ 是单态射，定义（\SGARefLink{sga7.sga7i.eix.n14-2-8}{14.2.8}）的这些扩张，恰好是来自正合列（\SGARefLink{sga7.sga7i.eix.n14-2-1}{14.2.1}）的扩张（\SGARefLink{sga7.sga7i.eix.n14-2-2}{14.2.2}）。这些扩张由纯代数方法定义，所以扩张（\SGARefLink{sga7.sga7i.eix.n14-2-8}{14.2.8}）也如此。事实上，鉴于 \SGARefLink{sga7.sga7i.eix.n11-11}{11.11} 将后一个群的连通分支群同构到 $\underline M_K\otimes\mathbb Q/\mathbb Z$，刚才给出的刻画立即表明，此扩张典范同构于 (8.2.3) 中引入、同样记为 $A_K^b$ 的群。

% END INCLUDED SGA7I_English_current_buildable_source_20260823/components/183_expose_IX_section_14_2.tex

% BEGIN INCLUDED SGA7I_English_current_buildable_source_20260823/components/184_expose_IX_section_14_3.tex
\par\medskip
\noindent
\SGARefTarget{sga7.sga7i.eix.n14-3}{14.3}. \underline{应用于下列同态的算术刻画：}
$i:\underline M_K\longrightarrow G_K$.

所考虑的刻画基于如下观察：对每个素数 $\ell$，由 $i$ 得到的 以 $T_\ell(G_K)$ 为核、以 $\underline M_K\otimes\mathbb Z_\ell$ 为商的扩张（\SGARefLink{sga7.sga7i.eix.n14-2-2}{14.2.2}）可由纯代数方法得知。特别地，$i$ 在典范同态

\begin{equation}\tag{\SGARefTarget{sga7.sga7i.eix.n14-3-1}{14.3.1}}
 \operatorname{Hom}(\underline M_K,G_K)\longrightarrow
 \prod_\ell\operatorname{Ext}^1
 \bigl(\underline M_K\otimes\mathbb Z_\ell,T_\ell(G_K)\bigr)
 \qquad .
\end{equation}

\noindent
下的像已知。在 $\underline M_K$ 为取值 $M$ 的常值格的情形（总可借助下降归结到此），上述同态可解释为

\begin{equation}\tag{\SGARefTarget{sga7.sga7i.eix.n14-3-2}{14.3.2}}
 \operatorname{Hom}\bigl(M,G_K(K)\bigr)\longrightarrow
 \prod_\ell\operatorname{Hom}
 \bigl(M,H^1(K,T_\ell(G_K))\bigr),
\end{equation}

\noindent
它由典范同态

\begin{equation}\tag{\SGARefTarget{sga7.sga7i.eix.n14-3-3}{14.3.3}}
 G_K(K)\longrightarrow\prod_\ell H^1(K,T_\ell(G_K))
 \qquad ,
\end{equation}

\noindent
得到，后者来自 Kummer 正合列中的连接同态

\begin{equation*}
 0\longrightarrow{}_nG_K\longrightarrow G_K
 \xrightarrow{\ n\ }G_K\longrightarrow0\qquad .
\end{equation*}

（事实上，同态（\SGARefLink{sga7.sga7i.eix.n14-3-1}{14.3.1}）总可解释为扭转群 $\check{\underline M}_K\otimes G_K=\underline{\operatorname{Hom}}(\underline M_K,G_K)$ 的、类似于（\SGARefLink{sga7.sga7i.eix.n14-3-3}{14.3.3}）的连接同态。）（\SGARefLink{sga7.sga7i.eix.n14-3-3}{14.3.3}）的核显然由 $G_K(K)$ 中\underline{对每个整数 $n\geq1$ 都可被 $n$ 除}的元素组成，我们称它们为\underline{无限可除}的元素。因此，这已经在模去 $\operatorname{Hom}(M,g)$ 的一个元素的意义下确定 $i$，其中 $g$ 是 $G_K(K)$ 的无限可除元素组成的子群：

\begin{equation*}
 g=\bigcap_{n\geq1}nG_K(K)\qquad .
\end{equation*}

现在有如下结果。

\medskip
\noindent
\underline{引理 \SGARefTarget{sga7.sga7i.eix.n14-3-4}{14.3.4}.} \underline{设 $V$ 是一个离散赋值环，其剩余域 $k$ 的特征为 $p>0$，且在其素域上为有限型；设 $G_K$ 是 $K$ 上的有限型群概形，若 $\operatorname{char}K=0$，则要求其幺幂秩为零。那么 $G_K(K)$ 的无限可除元素所成的群 $g$ 为零；即（若 $G_K$ 可除），同态}（\SGARefLink{sga7.sga7i.eix.n14-3-3}{14.3.3}）\underline{是单射。}

作有限扩张 $K'/K$，并以 $V$ 在 $K'$ 中的正规化代替它后，可假设 $G_K$ 有一个合成列，其相继商为以下类型之一：常值有限群、阿贝尔簇、$\mathbb G_{mK}$ 或 $\mathbb G_{aK}$。于是立即归结为 $G_K$ 本身为这些类型之一的情形。常值有限群或 $\mathbb G_{aK}$ 的情形平凡。在阿贝尔情形，可将 $G_K$ 看作其 N\'{e}ron 模型 $G$ 的泛纤维，故归结为说明：若 $G$ 是 $S$ 上的有限型概形，则 $G(S)$ 中无限可除的点组成的子群为零。

群 $G(k)$ 满足此结论，这由 $G_o$ 的通常逐层分析及\underline{Mordell--Weil 定理} [16] 立即得到（该定理说，若 $G_o$ 是有限型域 $k$ 上的阿贝尔概形，则 $G_o(k)$ 是有限生成 $\mathbb Z$-模）。因此 $g$ 包含于 $G(S)\longrightarrow G(k)$ 的核 $G(S)^{(1)}$。核 $G(S)^{(1)}\longrightarrow G(V/\underline m^{n+1})$ 有一个滤过，其商同构于 $k$ 上的向量群（此处已假设 $S$ 完备，这是允许的），故被 $p$ 零化。因此对每个 $n$，$g$ 都包含于前述同态的核，故它为零。$G_K=\mathbb G_{mK}$ 的情形归结于上述情形，因为 $G_K(K)=K^*$ 是 $\mathbb Z$（没有非零的无限可除元素）对 $V^*=\mathbb G_m(S)$ 的扩张，而刚才已证明后者满足同一条件。

\medskip
\noindent
\underline{注 \SGARefTarget{sga7.sga7i.eix.n14-3-4-1}{14.3.4.1}.} 注意，在 \SGARefLink{sga7.sga7i.eix.n14-3-4}{14.3.4} 中要求可被\underline{所有}整数 $\geq1$ 除，而不仅是可被固定素数 $\ell$ 的各次幂除，这是关键的；等价地，必须考虑（\SGARefLink{sga7.sga7i.eix.n14-3-3}{14.3.3}）的核，而不能将右端限制为单个因子 $H^1(K,T_\ell(G_K))$。事实上，若 $G_K$ 来自 $S$ 上的光滑群概形 $G$，且 $S$ 为 Hensel 的，则 $G(S)\longrightarrow G(k)$ 的核的每个元素，对每个 $\ell\neq p$ 都是无限 $\ell$-可除的！证明 \SGARefLink{sga7.sga7i.eix.n14-3-4}{14.3.4} 的论证表明，$G_K(K)$ 中无限 $p$-可除的元素组成的子群（即 $G_K$ 是 $p$-可除的时 $G_K(K)\longrightarrow H^1(K,T_p(G_K))$ 的核）是一个\underline{有限}群。

\medskip
\noindent
\SGARefTarget{sga7.sga7i.eix.n14-3-5}{14.3.5}. 上述讨论说明，当剩余域 $k$ 的特征为 $p>0$ 且在绝对意义下有限型时，同态 $i$（\SGARefLink{sga7.sga7i.eix.n14-1-5}{14.1.5}）完全由它在（\SGARefLink{sga7.sga7i.eix.n14-3-1}{14.3.1}）下的像，即由扩张（\SGARefLink{sga7.sga7i.eix.n14-2-2}{14.2.2}）确定。要从中得到一般情形下、对 $k$ 不作限制的 $i$ 的刻画原则，注意第7段中构造 $S$ 上概形 $G$ 的条件，远比 $S$ 为完备性状且其分式域上给定半稳定约化阿贝尔簇的条件更一般。

更一般地，设 $S$ 表示正规诺特环 $V$ 的谱，该环对某个 $\underline m$-进拓扑分离且完备，并考虑 $S$ 上一个纤维连通的光滑拟射影群概形 $A$。假设 $A_o=A\times_S S_o$ 是阿贝尔概形 $B_o$ 被等平凡环面 $T_o$ 的扩张，其中 $S_o=\operatorname{Spec}(V/\underline m)$；并假设 $A_U$ 是阿贝尔概形，其中 $U\subset S-S_o$ 为非空开子集。在7.22中构造了 $S$ 上的概形 $G$，它是阿贝尔概形 $B$ 被环面 $T$ 的扩张，由同构（\SGARefLink{sga7.sga7i.eix.n14-1-8}{14.1.8}）

\begin{equation}\tag{\SGARefTarget{sga7.sga7i.eix.n14-3-5-1}{14.3.5.1}}
 \hat G\xrightarrow{\ \sim\ }\hat A
\end{equation}

\noindent
在唯一同构意义下刻画，此同构定义在沿 $G_o$ 和 $A_o$ 的形式完备化上。对每个素数 $\ell$，得到同构（\SGARefLink{sga7.sga7i.eix.n7-3-5}{7.3.5}）

\begin{equation}\tag{\SGARefTarget{sga7.sga7i.eix.n14-3-5-2}{14.3.5.2}}
 T_\ell(G)\xrightarrow{\ \sim\ }T_\ell(A)^f\qquad ,
\end{equation}

\noindent
从而，限制到 $U$ 后得到 pro-$\ell$ Barsotti--Tate 群的扩张结构

\begin{equation}\tag{\SGARefTarget{sga7.sga7i.eix.n14-3-5-3}{14.3.5.3}}
\begin{tikzcd}[row sep=large,column sep=large]
0 \arrow[r] & T_\ell(G_U) \arrow[r] \arrow[d,equal]
 & T_\ell(A_U) \arrow[r]
 & T_\ell(A_U)/T_\ell(A_U)^f \arrow[r] & 0\quad .\\
& T_\ell(A_U)^f
\end{tikzcd}
\end{equation}

现在假设给定 $S$ 上的第二个群概形 $A'$，满足与 $A$ 相同的假设，连同 $(A,A')$ 对 $\mathbb G_m$ 的双扩张 $W$，它在 $U$ 上诱导阿贝尔概形 $A_U$ 和 $A'_U$ 之间的对偶性。（事实上，由 VIII \SGARefLink{sga7.sga7i.eviii.n7-1}{7.1}，给出 $W$ 就等价于给出 $A_U$ 和 $A'_U$ 间的这种对偶性。）于是 $W$ 定义配对

\begin{equation}\tag{\SGARefTarget{sga7.sga7i.eix.n14-3-5-4}{14.3.5.4}}
 T_\ell(A_U)\times T_\ell(A'_U)\longrightarrow\mathbb Z_\ell(1)
 \qquad ,
\end{equation}

\noindent
由此得到正交性定理（\SGARefLink{sga7.sga7i.eix.n7-4-8}{7.4.8}）：$T_\ell(A_U)^f$ 是 $T_\ell(A'_U)^t=T_\ell(T'_U)$ 的正交子群，其中 $T'$ 是与 $A'$ 联系的 Raynaud 扩张 $G'={A'}^{\natural}$ 的环面部分。因此再次得到同构

\begin{equation}\tag{\SGARefTarget{sga7.sga7i.eix.n14-3-5-5}{14.3.5.5}}
 T_\ell(A_U)/T_\ell(A_U)^f
 \simeq\underline M_U\otimes\mathbb Z_\ell\qquad ,
\end{equation}

\noindent
其中 $\underline M=D(T')$ 是 $S$ 上与环面 $T'$ 对偶的格。因此，（\SGARefLink{sga7.sga7i.eix.n14-3-5-3}{14.3.5.3}）和（\SGARefLink{sga7.sga7i.eix.n14-3-5-5}{14.3.5.5}）给出扩张

\begin{equation}\tag{\SGARefTarget{sga7.sga7i.eix.n14-3-5-6}{14.3.5.6}}
 0\longrightarrow T_\ell(G_U)\longrightarrow T_\ell(A_U)
 \longrightarrow\underline M_U\otimes\mathbb Z_\ell
 \longrightarrow0\qquad .
\end{equation}

在此理解下，可以陈述以下猜想。

% END INCLUDED SGA7I_English_current_buildable_source_20260823/components/184_expose_IX_section_14_3.tex

% BEGIN INCLUDED SGA7I_English_current_buildable_source_20260823/components/185_expose_IX_sections_14_4_and_14_5.tex
\par\medskip
\noindent
\underline{猜想 \SGARefTarget{sga7.sga7i.eix.n14-4}{14.4}.}  在刚才 \SGARefLink{sga7.sga7i.eix.n14-3-5}{14.3.5} 列出的假设下，存在同态

\begin{equation}\tag{\SGARefTarget{sga7.sga7i.eix.n14-4-1}{14.4.1}}
 i:\underline M_U\longrightarrow G_U
\end{equation}

\noindent
使对于每个素数 $\ell$，扩张（\SGARefLink{sga7.sga7i.eix.n14-3-5-6}{14.3.5.6}）同构于按（\SGARefLink{sga7.sga7i.eix.n14-2-2}{14.2.2}）所说明的方式由 $i$ 得到的扩张。更精确地说，把满足上述条件的同态（\SGARefLink{sga7.sga7i.eix.n14-4-1}{14.4.1}）联系于 \SGARefLink{sga7.sga7i.eix.n14-3-5}{14.3.5} 所考虑类型的每种情形，有且仅有一种方式，使 $i$ 的形成与每个由诺特进环的连续同态 $V\longrightarrow V'$ 所联系的基变换 $S'\longrightarrow S$ 相容。

\medskip
\noindent
\SGARefTarget{sga7.sga7i.eix.n14-4-2}{14.4.2}. 看来，与基变换相容的 $i$ 的存在性来自 M. Raynaud 的刚性解析理论，因为他是在比 \SGARefLink{sga7.sga7i.eix.n14-1}{14.1} 所述更一般的框架中发展此理论，像这里一样在一般诺特进环上工作。此外，为证明 \SGARefLink{sga7.sga7i.eix.n14-4}{14.4} 所断言的存在性与唯一性，可立即归结为只考虑 $S_o$ 为 $\mathbb Z$ 上有限型概形的那些 $S$。事实上，$V$ 是其在 $\mathbb Z$ 上有限型的子环 $V_i$ 的滤过归纳极限；由 Nagata 定理（EGA IV 7.8）可假设这些子环正规，而且对充分大的 $i$，情形 $(A,A',W)$ 来自 $V_i$ 上的类似情形，或者来自 $V_i$ 对由 $m_i=\underline m V_i$ 定义的拓扑的分离完备化上的类似情形。这就给出所述归结。

\medskip
\noindent
\SGARefTarget{sga7.sga7i.eix.n14-4-3}{14.4.3}. 现在注意 $S-S_o$ 是 Jacobson 概形（EGA IV 10.5.8）。因此，要说明给定情形 $(A,A',W)$ 的问题的两个解 $i,i'$ 相等，只须证明这些同态在 $U$ 的每个\underline{闭}点 $\eta$ 处相等。设 $S'$ 为 $\eta$ 在 $S$ 中的闭包的正规化。则 $S'$ 是完备性状，可将 \SGARefLink{sga7.sga7i.eix.n4-3}{4.3} 的唯一性结果应用于 $S$ 上情形到 $S'$ 的拉回，因为 $S'$ 的剩余域为\underline{有限域}。唯一性一经这样建立，在存在性的前提下，它就蕴含 $i$ 的形成与基变换相容。

\medskip
\noindent
\underline{注 \SGARefTarget{sga7.sga7i.eix.n14-5}{14.5}.} 还须以所需的一般性表述并建立以下相容性：\SGARefLink{sga7.sga7i.eix.n14-4}{14.4} 中用刚性解析方法引入的同态 $i$，与复解析理论中类似的同态（\SGARefLink{sga7.sga7i.eix.n13-5-13}{13.5.13}）之间的相容性——或者更好，找出两者的共同推广！


% END INCLUDED SGA7I_English_current_buildable_source_20260823/components/185_expose_IX_sections_14_4_and_14_5.tex


% BEGIN INCLUDED SGA7I_English_current_buildable_source_20260823/components/186_expose_IX_bibliography.tex
\bigskip
\noindent\underline{参考文献}

\begin{enumerate}
\item[{[1]}] M. Artin, \emph{Algebraization of formal moduli I}
  （待刊）。
\item[{[2]}] H. Cartan, \emph{Familles d'espaces complexes et fondements
  de la g\'{e}om\'{e}trie analytique}, S\'{e}minaire ENS 1960/61.
\item[{[3]}] P. Deligne, D. Mumford, \emph{The irreducibility of the space
  of curves of given genus}, Publ. Math. 第36期 (1969).
\item[{[4]}] P. Deligne, 正在撰写的 Hodge 理论文献（拟于1970年刊于 Publ. Math.）。
\item[{[5]}] M. Demazure, J. Giraud, M. Raynaud, \emph{Sch\'{e}mas
  ab\'{e}liens}, S\'{e}minaire Orsay 1967/68，待刊。
\item[{[6]}] A. Grothendieck, \emph{Le Groupe de Brauer III}, 载
  \emph{Dix Expos\'{e}s sur la cohomologie des sch\'{e}mas}, North-Holland
  Publishing Co.
\item[{[7]}] A. Grothendieck, \emph{Crystals and the De Rham cohomology of
  schemes} （I. Coates 和 O. Jussila 记录），载 \emph{Dix expos\'{e}s
  sur la cohomologie des sch\'{e}mas}, North-Holland Publishing Co.
\item[{[8]}] A. Grothendieck, \emph{Un th\'{e}or\`{e}me sur les
  homomorphismes de sch\'{e}mas ab\'{e}liens}, Invent. Math. \textbf{2}
  (1966), 第59--78页.
\item[{[9]}] P. A. Griffiths, \emph{Report on variation of Hodge
  structures} （待刊）。
\item[{[10]}] M. Hakim, \emph{Topos annel\'{e}s et sch\'{e}mas relatifs},
  拟刊于 Grundlehren der Mathematik，Springer（1970）。
\item[{[11]}] J. I. Igusa, \emph{Abstract vanishing cycle theory},
  Proceedings of the Japan Academy \textbf{34}, 第9期 (1958),
  第589--593页.
\item[{[12]}] R. Kiehl, \emph{Der Endlichkeitssatz f\"{u}r eigentliche
  Abbildungen}, Invent. Math. \textbf{2} (1967), 第191--214页.
\item[{[13]}] R. Kiehl.
\item[{[14]}] S. L. Kleiman, \emph{Algebraic cycles and the Weil
  conjectures}, in \emph{Dix expos\'{e}s sur la cohomologie des sch\'{e}mas},
  North-Holland Publishing Co.
\item[{[15]}] S. Lang, \emph{Abelian varieties}, Interscience Publishers,
  第7期 (1959).
\item[{[16]}] S. Lang and A. N\'{e}ron, \emph{Rational points of abelian
  varieties over function fields}, Amer. J. Math. \textbf{81}, 第1期
  (1959), 第95--118页.
\item[{[16bis]}] U. I. Manin, \emph{Th\'{e}orie des groupes formels
  commutatifs}, Uspekhi Mat. Nauk \textbf{XVIII} (1963), 第4--90页.
\item[{[17]}] D. Mumford, \emph{Geometric invariant theory}, Ergebnisse
  der Mathematik 34, Springer (1965).
\item[{[18]}] D. Mumford, \emph{Abstract Theta functions}, Advanced
  Science Seminar in Algebraic Geometry, Bowdoin College, 1967 （H. Pittie 记录）。
\item[{[19]}] A. N\'{e}ron, \emph{Mod\`{e}les minimaux des vari\'{e}t\'{e}s
  ab\'{e}liennes sur les corps locaux et globaux}, Publ. Math. 21 (1964),
  第5--128页.
\item[{[20]}] M. Raynaud, \emph{Caract\'{e}ristique d'Euler--Poincar\'{e}
  d'un faisceau et cohomologie des vari\'{e}t\'{e}s ab\'{e}liennes}, 载
  \emph{Dix expos\'{e}s sur la cohomologie des sch\'{e}mas}, North-Holland
  Publishing Co.
\item[{[21]}] M. Raynaud, \emph{Mod\`{e}les de N\'{e}ron}, C. R. Acad.
  Sci. Paris \textbf{262} （1966年2月）, 第413--416页.
\item[{[22]}] M. Raynaud, \emph{Sp\'{e}cialisation du foncteur de Picard},
  C. R. Acad. Sci. \textbf{264}, 第941--943和1001--1004页。
\item[{[22bis]}] M. Raynaud, \emph{Sp\'{e}cialisation du foncteur de
  Picard}, 拟刊于 Publ. Math.
\item[{[23]}] M. Raynaud, \emph{Faisceaux amples sur les sch\'{e}mas en
  groupes et les espaces homog\`{e}nes}, Paris 学位论文，1968（拟于1970年刊于 Lecture Notes，Springer）。
\item[{[24]}] M. Raynaud.
\item[{[25]}] M. Rosenlicht, \emph{Some basic theorems on algebraic
  groups}, Amer. J. Math. \textbf{78}, 第2期 (1956), 第401--443页.
\item[{[25bis]}] J.-P. Serre, \emph{Quelques propri\'{e}t\'{e}s des
  vari\'{e}t\'{e}s ab\'{e}liennes en caract\'{e}ristique $p$}, Amer. J.
  Math. \textbf{80} (1958), 第715--739页.
\item[{[26]}] J.-P. Serre, \emph{Corps locaux}, Act. Sci. Ind. 1296
  (1962), Hermann, Paris.
\item[{[27]}] J.-P. Serre, \emph{Sur les corps locaux \`{a} corps
  r\'{e}siduel alg\'{e}briquement clos}, Bull. Soc. Math. France
  \textbf{89} (1961), 第105--154页.
\item[{[28]}] J.-P. Serre, \emph{Abelian $\ell$-adic representations and
  elliptic curves}, Benjamin Inc. (1968).
\item[{[29]}] J.-P. Serre and J. Tate, \emph{Good reduction of abelian
  varieties}, Ann. of Math. \textbf{88}, 第3期 (1968), 第492--517页.
\item[{[30]}] J.-P. Serre and J. Tate, \emph{Elliptic curves and formal
  groups} （代数几何暑期研究班讲义，Woods Hole，1964）。
\item[{[31]}] J.-P. Serre, \emph{Morphismes universels et vari\'{e}t\'{e}
  d'Albanese}, S\'{e}minaire Chevalley 1958/59，第10讲。
\item[{[32]}] J. Tate, \emph{WC-groups over $p$-adic fields}, S\'{e}minaire
  Bourbaki 第156期 (1957).
\item[{[33]}] J. Tate, \emph{$p$-divisible groups}, Proceedings of a
  Conference on Local Fields, Driebergen (1966), Springer.
\item[{[34]}] J. Tate, \emph{Rigid analytic spaces} （IHES 油印）。
\end{enumerate}


% END INCLUDED SGA7I_English_current_buildable_source_20260823/components/186_expose_IX_bibliography.tex

% BEGIN INCLUDED SGA7I_English_current_buildable_source_20260823/components/187_expose_IX_terminal_publisher_matter.tex
% R3 explicit revisions IX-B-14; see qa/FIDELITY-R3-CHANGES.json.
\clearpage
\begingroup
\small
\sloppy
Lecture Notes in Mathematics


详细书目可函索

第111卷： K. H. Mayer, Relationen zwischen charakteristischen Zahlen. III, 99页. 1969. DM 16,-

第112卷： Colloquium on Methods of Optimization. 编者：N. N. Moiseev. IV, 293页. 1970. DM 18,-

第113卷： R. Wille, Kongruenzklassengeometrien. III, 99页. 1970. DM 16,-

第114卷： H. Jacquet and R. P. Langlands, Automorphic Forms on GL (2). VII, 548页. 1970. DM 24,-

第115卷： K. H. Roggenkamp and V. Huber-Dyson, Lattices over Orders I. XIX, 290页. 1970. DM 18,-

第116卷： Sèminaire Pierre Lelong (Analyse) Année 1969. IV, 195页. 1970. DM 16,-

第117 卷： Y. Meyer, Nombres de Pisot, Nombres de Salem et Analyse Harmonique. 63页. 1970. DM 16,-

第118卷： Proceedings of the 15th Scandinavian Congress, Oslo 1968. 编者：K. E. Aubert and W. Ljunggren. IV, 162页. 1970. DM 16,-

第119卷： M. Raynaud, Faisceaux amples sur les schémas en groupes et les espaces homogènes. III, 219页. 1970. DM 16,-

第120卷： D. Siefkes, Büchi's Monadic Second Order Successor Arithmetic. XII, 130页. 1970. DM 16,-

第121卷： H. S. Bear, Lectures on Gleason Parts. III, 47页. 1970. DM 16,-

第122卷： H. Zieschang, E. Vogt und H.-D. Coldewey, Flächen und ebene diskontinuierliche Gruppen. VIII, 203页. 1970. DM 16,-

第123卷： A. V. Jategaonkar, Left Principal Ideal Rings.VI, 145页. 1970. DM 16,-

第124卷： Séminare de Probabilités IV. 编者：P. A. Meyer. IV, 282页. 1970. DM 20,-

第125卷： Symposium on Automatic Demonstration. V, 310页.1970. DM 20,-

第126卷： P. Schapira, Théorie des Hyperfonctions. XI,157页.1970. DM 16,-

第127卷： I. Stewart, Lie Algebras. IV, 97页. 1970. DM 16,-

第128卷： M. Takesaki, Tomita's Theory of Modular Hilbert Algebras and its Applications. II, 123页. 1970. DM 16,-

第129卷： K. H. Hofmann, The Duality of Compact Semigroups and C$^*$- Bigebras. XII, 142页. 1970. DM 16,-

第130卷： F. Lorenz, Quadratische Formen über Körpern. II, 77页. 1970. DM 16,-

第131卷： A Borel et al., Seminar on Algebraic Groups and Related Finite Groups. VII, 321页. 1970. DM 22,-

第132卷： Symposium on Optimization. III, 348页. 1970. DM 22,-

第133卷： F. Topsøe, Topology and Measure. XIV, 79页. 1970. DM 16,-

第134卷： L. Smith, Lectures on the Eilenberg-Moore Spectral Sequence. VII, 142页. 1970. DM 16,-

第135卷： W. Stoll, Value Distribution of Holomorphic Maps into Compact Complex Manifolds. II, 267页. 1970. DM 18,-

第136 卷： M. Karoubi et al., Séminaire Heidelberg-Saarbrücken-Strasbuorg sur la K-Théorie. IV, 264页. 1970. DM 18,-

第137 卷： Reports of the Midwest Category Seminar IV. 编者：S. MacLane. III, 139页. 1970. DM 16,-

第138卷： D. Foata et M. Schützenberger, Théorie Géométrique des Polynômes Eulériens. V, 94页. 1970. DM 16,-

第139卷： A. Badrikian, Séminaire sur les Fonctions Aléatoires Linéaires et les Mesures Cylindriques. VII, 221页. 1970. DM 18,-

第140卷： Lectures in Modern Analysis and Applications II. 编者：C. T. Taam. VI, 119页. 1970. DM 16,-

第141卷： G. Jameson, Ordered Linear Spaces. XV, 194页. 1970. DM 16,-

第142卷： K. W. Roggenkamp, Lattices over Orders II. V, 388页. 1970. DM 22,-

第143卷： K. W. Gruenberg, Cohomological Topics in Group Theory. XIV, 275页. 1970. DM 20,-

第144卷： Seminar on Differential Equations and Dynamical Systems, II. 编者：J. A. Yorke. VIII, 268页. 1970. DM 20,-

第145卷： E. J. Dubuc, Kan Extensions in Enriched Category Theory. XVI, 173页. 1970. DM 16,-

第146卷： A. B. Altman and S. Kleiman, Introduction to Grothendieck Duality Theory. II, 192页. 1970. DM 18,-

第147卷： D. E. Dobbs, Cech Cohomological Dimensions for Commutative Rings. VI, 176页. 1970. DM 16,-

第148卷： R. Azencott, Espaces de Poisson des Groupes Localement Compacts. IX, 141页. 1970. DM 16,-

第149卷： R. G. Swan and E. G. Evans, K-Theory of Finite Groups and Orders. IV, 237页. 1970. DM 20,-

第150卷： Heyer, Dualität lokalkompakter Gruppen. XIII, 372页. 1970. DM 20,-

第151卷： M. Demazure et A. Grothendieck, Schémas en Groupes I. (SGA 3). XV, 562页. 1970. DM 24,-

第152卷： M. Demazure et A. Grothendieck, Schémas en Groupes II. (SGA 3). IX, 654页. 1970. DM 24,-

第153卷： M. Demazure et A. Grothendieck, Schémas en Groupes III. (SGA 3). VIII, 529页. 1970. DM 24,-

第154卷： A. Lascoux et M. Berger, Variétés Kähleriennes Compactes. VII, 83页. 1970. DM 16,-

第155卷： Several Complex Variables I, Maryland 1970. 编者：J. Horváth. IV, 214页. 1970. DM 18,-

第156卷： R. Hartshorne, Ample Subvarieties of Algebraic Varieties. XIV, 256页. 1970. DM 20,-

第157卷： T. tom Dieck, K. H. Kamps und D. Puppe, Homotopietheorie. VI, 265页. 1970. DM 20,-

第158卷： T. G. Ostrom, Finite Translation Planes. IV. 112页. 1970. DM 16,-

第159卷： R. Ansorge und R. Hass. Konvergenz von Differenzenverfahren für lineare und nichtlineare Anfangswertaufgaben. VIII, 145页. 1970. DM 16,-

第160卷： L. Sucheston, Constributions to Ergodic Theory and Probability. VII, 277页. 1970. DM 20,-

第161卷： J. Stasheff, H-Spaces from a Homotopy Point of View. VI, 95页. 1970. DM 16,-

第162卷： Harish-Chandra and van Dijk, Harmonic Analysis on Reductive p-adic Groups. IV, 125页. 1970. DM 16,-

第163卷： P. Deligne, Equations Différentielles à Points Singuliers Reguliers. III, 133页. 1970. DM 16,-

第164卷： J. P. Ferrier, Seminaire sur les Algebres Complètes. II, 69页. 1970. DM 16,-

第165卷： J. M. Cohen, Stable Homotopy. V, 194页. 1970. DM 16,-

第166卷： A. J. Silberger, PGL$_2$ over the p-adics: its Representations, Spherical Functions, and Fourier Analysis. VII, 202页. 1970. DM 18,-

第167卷： Lavrentiev, Romanov and Vasiliev, Multidimensional Inverse Problems for Differential Equations. V, 59页. 1970. DM 16,-

第168卷： F. P. Peterson, The Steenrod Algebra and its Applications: A conference to Celebrate N. E. Steenrod's Sixtieth Birthday. VII, 317页. 1970. DM 22,-

第169卷： M. Raynaud, Anneaux Locaux Henséliens. V, 129页. 1970. DM 16,-

第170卷： Lectures in Modern Analysis and Applications III. 编者：C. T. Taam. VI, 213页. 1970. DM 18,-

第171卷： Set-Valued Mappings, Selections and Topological Properties of $2^X$. 编者：W. M. Fleischman. X, 110页. 1970. DM 16,-

第172卷： Y.-T. Siu and G. Trautmann, Gap-Sheaves and Extension of Coherent Analytic Subsheaves. V, 172页. 1971. DM 16,-

第173卷： J. N. Mordeson and B. Vinograde, Structure of Arbitrary Purely Inseparable Extension Fields. IV, 138页. 1970. DM 16,-

第174卷： B. Iversen, Linear Determinants with Applications to the Picard Scheme of a Family of Algebraic Curves. VI, 69页. 1970. DM 16,-

第175卷： M. Brelot, On Topologies and Boundaries in Potential Theory. VI, 176页. 1971. DM 18,-

第176卷： H. Popp, Fundamentalgruppen algebraischer Mannigfaltigkeiten. IV, 154页. 1970. DM 16,-

第177卷： J. Lambek, Torsion Theories, Additive Semantics and Rings of Quotients. VI, 94页. 1971. DM 16,-

请翻页


%%%%%%%%%%  scan idx 537  |  folio 526  %%%%%%%%%%

第178卷： Th. Bröcker und T. tom Dieck, Kobordismentheorie. XVI, 191页. 1970. DM 18,-

第179卷： Seminaire Bourbaki -- vol. 1968/69. Exposés 347-363. IV. 295页. 1971. DM 22,-

第180卷： Séminaire Bourbaki -- vol. 1969/70. Exposés 364-381. IV, 310页. 1971. DM 22,-

第181卷： F. DeMeyer and E. Ingraham, Separable Algebras over Commutative Rings. V, 157页. 1971. DM 16.-

第182卷： L. D. Baumert. Cyclic Difference Sets. VI, 166页. 1971. DM 16,-

第183卷： Analytic Theory of Differential Equations. 编者：P. F. Hsieh and A. W. J. Stoddart. VI, 225页. 1971. DM 20,-

第184卷： Symposium on Several Complex Variables, Park City, Utah, 1970. 编者：R. M. Brooks. V, 234页. 1971. DM 20,-

第185卷： Several Complex Variables II, Maryland 1970. 编者：J. Horváth. III, 287页. 1971. DM 24,-

第186卷： Recent Trends in Graph Theory. 编者：M. Capobianco/J. B. Frechen/M. Krolik. VI, 219页. 1971. DM 18.-

第187卷： H. S. Shapiro, Topics in Approximation Theory. VIII, 275页. 1971. DM 22,-

第188卷： Symposium on Semantics of Algorithmic Languages. 编者：E. Engeler. VI, 372页. 1971. DM 26,-

第189卷： A. Weil, Dirichlet Series and Automorphic Forms. V, 164页. 1971. DM 16,-

第190卷： Martingales. A Report on a Meeting at Oberwolfach, May 17-23, 1970. 编者：H. Dinges. V, 75页. 1971. DM 16,-

第191卷： Séminaire de Probabilités V. 编者：P. A. Meyer. IV, 372页. 1971. DM 26,-

第192卷： Proceedings of Liverpool Singularities -- Symposium I. 编者：C. T. C. Wall. V, 319页. 1971. DM 24,-

第193卷： Symposium on the Theory of Numerical Analysis. 编者：J. Ll. Morris. VI, 152页. 1971. DM 16,-

第194卷： M. Berger, P. Gauduchon et E. Mazet. Le Spectre d'une Variété Riemannienne. VII, 251页. 1971. DM 22,-

第195卷： Reports of the Midwest Category Seminar V. 编者：J.W. Gray and S. Mac Lane.III, 255页. 1971. DM 22,-

第196卷： H-spaces -- Neuchâtel (Suisse)- Août 1970. 编者：F. Sigrist, V, 156页. 1971. DM 16,-

第197卷： Manifolds -- Amsterdam 1970. 编者：N. H. Kuiper. V, 231页. 1971. DM 20,-

第198卷： M. Hervé, Analytic and Plurisubharmonic Functions in Finite and Infinite Dimensional Spaces. VI, 90页. 1971. DM 16.-

第199卷： Ch. J. Mozzochi, On the Pointwise Convergence of Fourier Series. VII, 87页. 1971. DM 16,-

第200卷： U. Neri, Singular Integrals. VII, 272页. 1971. DM 22,-

第201 卷： J. H. van Lint, Coding Theory. VII, 136页. 1971. DM 16,-

第202卷： J. Benedetto, Harmonic Analysis on Totally Disconnected Sets. VIII, 261页. 1971. DM 22,-

第203卷： D. Knutson, Algebraic Spaces. VI, 261页. 1971. DM 22,-

第204卷： A. Zygmund, Intégrales Singulières. IV, 53页. 1971. DM 16,-

第205卷： Séminaire Pierre Lelong (Analyse) Année 1970. VI, 243页. 1971. DM 20,-

第206卷： Symposium on Differential Equations and Dynamical Systems. 编者：D. Chillingworth. XI, 173页. 1971. DM 16,-

第207卷： L. Bernstein, The Jacobi-Perron Algorithm -- Its Theory and Application. IV, 161页. 1971. DM 16,-

第208卷： A. Grothendieck and J. P. Murre, The Tame Fundamental Group of a Formal Neighbourhood of a Divisor with Normal Crossings on a Scheme. VIII, 133页. 1971. DM 16,-

第209卷： Proceedings of Liverpool Singularities \quad Symposium II. 编者：C. T. C. Wall. V, 280页. 1971. DM 22,-

第210卷： M. Eichler, Projective Varieties and Modular Forms. III, 118页. 1971. DM 16,-

第211卷： Théorie des Matroïdes. 编者：C. P. Bruter. III, 108页. 1971. DM 16,-

第212卷： B. Scarpellini, Proof Theory and Intuitionistic Systems. VII, 291页. 1971. DM 24,-

第213卷： H. Hogbe-Nlend, Théorie des Bornologies et Applications. V, 168页. 1971. DM 18,-

第214卷： M. Smorodinsky, Ergodic Theory, Entropy. V, 64页. 1971. DM 16,-

第215卷： P. Antonelli, D. Burghelea and P. J. Kahn, The Concordance-Homotopy Groups of Geometric Automorphism Groups. X, 140页. 1971. DM 16,-

第216卷： H. Maaß, Siegel's Modular Forms and Dirichlet Series. VII, 328页. 1971. DM 20,-

第217卷： T. J. Jech, Lectures in Set Theory with Particular Emphasis on the Method of Forcing. V, 137页. 1971. DM 16,-

第218卷： C. P. Schnorr, Zufälligkeit und Wahrscheinlichkeit. IV, 212页 1971. DM 20,-

第219卷： N. L. Alling and N. Greenleaf, Foundations of the Theory of Klein Surfaces. IX, 117页. 1971. DM 16,-

第220卷： W. A. Coppel, Disconjugacy. V, 148页. 1971. DM 16,-

第221卷： P. Gabriel und F. Ulmer, Lokal präsentierbare Kategorien. V, 200页. 1971. DM 18,-

第222卷： C. Meghea, Compactification des Espaces Harmoniques. III, 108页. 1971. DM 16,-

第223卷： U. Felgner, Models of ZF-Set Theory. VI, 173页. 1971. DM 16,-

第224卷： Revètements Etales et Groupe Fondamental. (SGA 1). 主持：A. Grothendieck XXII, 447页. 1971. DM 30,-

第225卷： Théorie des Intersections et Théorème de Riemann-Roch. (SGA 6). 主持：P. Berthelot, A. Grothendieck et L. Illusie. XII, 700页. 1971. DM 40,-

第226卷： Seminar on Potential Theory, II. 编者：H. Bauer. IV, 170页. 1971. DM 18,-

第227卷： H. L. Montgomery, Topics in Multiplicative Number Theory. IX, 178页. 1971. DM 18,-

第228卷： Conference on Applications of Numerical Analysis. 编者：J. Ll. Morris. X, 358页. 1971. DM 26,-

第229卷： J. Väisälä, Lectures on n-Dimensional Quasiconformal Mappings. XIV, 144页. 1971. DM 16,-

第230卷： L. Waelbroeck, Topological Vector Spaces and Algebras. VII, 158页. 1971. DM 16,-

第231卷： H. Reiter, L$^1$-Algebras and Segal Algebras. XI, 113页. 1971. DM 16,-

第232卷： T. H. Ganelius, Tauberian Remainder Theorems. VI, 75页. 1971. DM 16,-

第233卷： C. P. Tsokos and W. J. Padgett. Random Integral Equations with Applications to Stochastic Systems. VII, 174页. 1971. DM 18,-

第234卷： A. Andreotti and W. Stoll. Analytic and Algebraic Dependence of Meromorphic Functions. III, 390页. 1971. DM 26,-

第235卷： Global Differentiable Dynamics. 编者：O. Hájek, A. J. Lohwater, and R. McCann. X, 140页. 1971. DM 16,-

第236 卷： M. Barr, P. A. Grillet, and D. H. van Osdol. Exact Categories and Categories of Sheaves. VII, 239页. 1971, DM 20,-

第237卷： B. Stenström. Rings and Modules of Quotients. VII, 136页. 1971. DM 16,-

第238卷： Der kanonische Modul eines Cohen-Macaulay-Rings. 编者：Jürgen Herzog und Ernst Kunz. VI, 103页. 1971. DM 16,-

第239卷： L. Illusie, Complexe Cotangent et Déformations I. XV, 355页. 1971. DM 26,-

第240卷： A. Kerber, Representations of Permutation Groups I. VII, 192页. 1971. DM 18,-

第241卷： S. Kaneyuki, Homogeneous Bounded Domains and Siegel Domains. V, 89页. 1971. DM 16,-

第242卷： R. R. Coifman et G. Weiss, Analyse Harmonique Non-Commutative sur Certains Espaces. V, 160页. 1971. DM 16,-

第243卷： Japan-United States Seminar on Ordinary Differential and Functional Equations. 编者：M. Urabe. VIII, 332页. 1971. DM 26,-

第244卷： Séminaire Bourbaki -- vol. 1970/71. Exposés 382-399. IV, 356页. 1971. DM 26,-

第245卷： D. E. Cohen, Groups of Cohomological Dimension One. V, 99页. 1972. DM 16,-


%%%%%%%%%%  scan idx 538  |  folio 527  %%%%%%%%%%

\noindent\rule{\linewidth}{0.4pt}

第246卷： Lectures on Rings and Modules. Tulane University Ring and Operator Theory Year, 1970--1971. Volume I. X, 661页. 1972. DM 40,--

第247卷： Lectures on Operator Algebras. Tulane University Ring and Operator Theory Year, 1970--1971. Volume II. XI, 786页. 1972. DM 40,--

第248卷： Lectures on the Applications of Sheaves to Ring Theory. Tulane University Ring and Operator Theory Year, 1970--1971. Volume III. VIII, 315页. 1971. DM 26,--

第249卷： Symposium on Algebraic Topology. 编者：P. J. Hilton. VII, 111页. 1971. DM 16,--

第250卷： B. Jónsson, Topics in Universal Algebra. VI, 220页. 1972. DM 20,--

第251卷： The Theory of Arithmetic Functions. 编者：A. A. Gioia and D. L. Goldsmith VI, 287页. 1972. DM 24,--

第252卷： D. A. Stone, Stratified Polyhedra. IX, 193页. 1972. DM 18,--

第253卷： V. Komkov, Optimal Control Theory for the Damping of Vibrations of Simple Elastic Systems. V, 240页. 1972. DM 20,--

第254卷： C. U. Jensen, Les Foncteurs Dérivés de $\varprojlim$ et leurs Applications en Théorie des Modules. V, 103页. 1972. DM 16,--

第255卷： Conference in Mathematical Logic -- London '70. 编者：W. Hodges. VIII, 351页. 1972. DM 26,--

第256卷： C. A. Berenstein and M. A. Dostal, Analytically Uniform Spaces and their Applications to Convolution Equations. VII, 130页. 1972. DM 16,--

第257卷： R. B. Holmes, A Course on Optimization and Best Approximation. VIII, 233页. 1972. DM 20,--

第258卷： Séminaire de Probabilités VI. 编者：P. A. Meyer. VI, 253页. 1972. DM 22,--

第259卷： N. Moulis, Structures de Fredholm sur les Variétés Hilbertiennes. V, 123页. 1972. DM 16,--

第260卷： R. Godement and H. Jacquet, Zeta Functions of Simple Algebras. IX, 188页. 1972. DM 18,--

第261卷： A. Guichardet, Symmetric Hilbert Spaces and Related Topics. V, 197页. 1972. DM 18,--

第262卷： H. G. Zimmer, Computational Problems, Methods, and Results in Algebraic Number Theory. V, 103页. 1972. DM 16,--

第263卷： T. Parthasarathy, Selection Theorems and their Applications. VII, 101页. 1972. DM 16,--

第264卷： W. Messing, The Crystals Associated to Barsotti-Tate Groups: with Applications to Abelian Schemes. III, 190页. 1972. DM 18,--

第265卷： N. Saavedra Rivano, Catégories Tannakiennes. II, 418页. 1972. DM 26,--

第266卷： Conference on Harmonic Analysis. 编者：D. Gulick and R. L. Lipsman. VI, 323页. 1972. DM 24,--

第267卷： Numerische Lösung nichtlinearer partieller Differential- und Integro-Differentialgleichungen. 编者：R. Ansorge und W. Törnig, VI, 339页. 1972. DM 26,--

第268卷： C. G. Simader, On Dirichlet's Boundary Value Problem. IV, 238页. 1972. DM 20,--

第269卷： Théorie des Topos et Cohomologie Etale des Schémas. (SGA 4). 主持：M. Artin, A. Grothendieck et J. L. Verdier. XIX, 525页. 1972. DM 50,--

第270卷： Théorie des Topos et Cohomologie Etale des Schémas. Tome 2. (SGA 4). 主持：M. Artin, A. Grothendieck et J. L. Verdier. V, 314页. 1972. DM 50,--

第271卷： J. P. May, The Geometry of Iterated Loop Spaces. IX, 175页. 1972. DM 18,--

第272卷： K. R. Parthasarathy and K. Schmidt, Positive Definite Kernels, Continuous Tensor Products, and Central Limit Theorems of Probability Theory. VI, 107页. 1972. DM 16,--

第273卷： U. Seip, Kompakt erzeugte Vektorräume und Analysis. IX, 119页. 1972. DM 16,--

第274卷： Toposes, Algebraic Geometry and Logic. 编者：F. W. Lawvere. VI, 189页. 1972. DM 18,--

第275卷： Séminaire Pierre Lelong (Analyse) Année 1970--1971. VI, 181页. 1972. DM 18,--

第276卷： A. Borel, Représentations de Groupes Localement Compacts. V, 98页. 1972. DM 16,--

第277卷： Séminaire Banach. 编者：C. Houzel. VII, 229页. 1972. DM 20,--

第278卷： H. Jacquet, Automorphic Forms on GL(2). Part II. XIII, 142页. 1972. DM 16,--

第279卷： R. Bott, S. Gitler and I. M. James, Lectures on Algebraic and Differential Topology. V, 174页. 1972. DM 18,--

第280卷： Conference on the Theory of Ordinary and Partial Differential Equations. 编者：W. N. Everitt and B. D. Sleeman. XV, 367页. 1972. DM 26,--

第281卷： Coherence in Categories. 编者：S. Mac Lane. VII, 235页. 1972. DM 20,--

第282卷： W. Klingenberg und P. Flaschel, Riemannsche Hilbertmannigfaltigkeiten. Periodische Geodätische. VII, 211页. 1972. DM 20,--

第283卷： L Illusie, Complexe Cotangent et Déformations II. VII, 304页. 1972. DM 24,--

第284卷： P. A. Meyer, Martingales and Stochastic Integrals I. VI, 89页. 1972. DM 16,--

第285卷： P. de la Harpe, Classical Banach-Lie Algebras and Banach-Lie Groups of Operators in Hilbert Space. III, 160页. 1972. DM 16,--

第286卷： S. Murakami, On Automorphisms of Siegel Domains. V, 95页. 1972. DM 16,--

第288卷： Groupes de Monodromie en Géométrie Algébrique. (SGA 7 I). 主持：A. Grothendieck. IX, 523页. 1972. DM 50,--

\noindent\rule{\linewidth}{0.4pt}


%%%%%%%%%%  scan idx 539  |  folio 528  %%%%%%%%%%

ISBN 3-540-05987-3

ISBN 0-387-05987-3
\endgroup


% END INCLUDED SGA7I_English_current_buildable_source_20260823/components/187_expose_IX_terminal_publisher_matter.tex

